% Rozdzial 12 w kolejnosci cpp.tex — metaprogramowanie w C++20.
% Fragment do wlaczenia po rozdziale 11.
% Wymagane pakiety: listings, array, amsmath, hyperref.
% Uzycie: % Rozdzial 12 w kolejnosci cpp.tex — metaprogramowanie w C++20.
% Fragment do wlaczenia po rozdziale 11.
% Wymagane pakiety: listings, array, amsmath, hyperref.
% Uzycie: % Rozdzial 12 w kolejnosci cpp.tex — metaprogramowanie w C++20.
% Fragment do wlaczenia po rozdziale 11.
% Wymagane pakiety: listings, array, amsmath, hyperref.
% Uzycie: % Rozdzial 12 w kolejnosci cpp.tex — metaprogramowanie w C++20.
% Fragment do wlaczenia po rozdziale 11.
% Wymagane pakiety: listings, array, amsmath, hyperref.
% Uzycie: \input{rozdzial_9_metaprogramowanie_cpp.tex}

\chapter[Metaprogramowanie w C++]{Zaawansowane metaprogramowanie: szablony, cechy typów i obliczenia podczas kompilacji}
\label{chap:meta-cpp}

\section{Czym jest metaprogramowanie?}

\emph{Metaprogramowanie} oznacza pisanie kodu, który na etapie kompilacji
wyznacza typy, własności typów, wartości lub strukturę później
kompilowanego programu. Szablon nie jest gotową funkcją czy klasą:
stanowi przepis, z którego kompilator tworzy konkretne wersje
dla podanych typów. Nazywamy to \emph{instancjowaniem szablonu}.
Na przykład wywołanie
\texttt{twice(21)} w programie poniżej każe kompilatorowi przygotować
wersję funkcji dla \texttt{int}; program podczas działania wywołuje
już tę wersję. W C++20 do dyspozycji mamy kilka narzędzi:
\begin{itemize}
\item \textbf{szablony klas i funkcji} — generowanie kodu dla różnych typów;
\item \textbf{\texttt{type\_traits}} — obliczanie własności i przekształcanie typów;
\item \textbf{specjalizacje i SFINAE} — wybór poprawnych wersji szablonów;
\item \textbf{\texttt{if constexpr}, \texttt{constexpr}, \texttt{consteval}} —
      wykonywanie obliczeń i selekcja kodu w czasie kompilacji;
\item \textbf{\texttt{concepts}} — czytelne formułowanie wymagań wobec typów;
\item \textbf{CRTP} — wykorzystanie dziedziczenia do polimorfizmu statycznego.
\end{itemize}

Metaprogramowanie ma praktyczne cele: eliminowanie powtarzalnego kodu,
uzyskiwanie lepszych błędów kompilacji, wybór wydajnej implementacji oraz
zapewnienie poprawności określonych operacji \emph{przed uruchomieniem} programu.
Nie oznacza automatycznie zwiększenia wydajności — kosztem mogą być dłuższa
kompilacja, powiększenie kodu maszynowego i trudniejsze komunikaty diagnostyczne.

\begin{lstlisting}
#include <type_traits>

template<class T>
T twice(T x) {
    static_assert(std::is_arithmetic_v<T>);
    return x + x;
}

int a = twice(21);       // T = int
// twice("text");       // static_assert fails
\end{lstlisting}

\paragraph{Związki z innymi rozdziałami.}
Wykorzystamy szablony wariadyczne (rozdział~11),
przenoszenie i referencje (rozdział~6) oraz lambdy (rozdział~9).
Koncepcje szerzej poznamy w rozdziale~13; tutaj pokazujemy
przede wszystkim mechanizmy, na których się opierają.
Wszystkie oznaczone jako poprawne przykłady odnoszą się do \textbf{C++20}.

\section{Instancjowanie szablonów, parametry i zależne nazwy}
\label{sec:cpp9-instantiation}

\subsection{Parametry typowe, nietypowe i szablonowe}

Szablon może przyjmować argumenty będące typami, wartościami stałymi
lub innymi szablonami. Parametr wartości (\textit{non-type template parameter},
NTTP) w C++20 może również mieć odpowiedni typ klasowy spełniający
wymagania tzw. typu strukturalnego.

\begin{lstlisting}
#include <array>
#include <cstddef>

template<class T, std::size_t N>
struct Buffer {
    std::array<T, N> data{};
};

Buffer<double, 4> samples; // T=double, N=4
\end{lstlisting}

\texttt{Buffer<double,4>} i \texttt{Buffer<double,8>} są różnymi typami.
Parametr \texttt{N} jest znany podczas kompilacji i może wystąpić tam,
gdzie język wymaga wyrażenia stałego, np. jako rozmiar \texttt{std::array}.

\subsection{Parametr będący obiektem klasy (C++20)}

Od C++20 odpowiednie \emph{typy strukturalne} mogą występować jako
parametry szablonu będące wartościami. Umożliwia to np. identyfikowanie
szablonu przez napis znany podczas kompilacji:

\begin{lstlisting}
#include <cstddef>

template<std::size_t N>
struct fixed_string {
    char data[N]{};
    constexpr fixed_string(const char (&s)[N]) {
        for (std::size_t i = 0; i < N; ++i)
            data[i] = s[i];
    }
};

template<fixed_string Name>
struct NamedCommand {
    static constexpr auto name = Name;
};

static_assert(NamedCommand<"run">::name.data[0] == 'r');
\end{lstlisting}

To przykład zaawansowanej techniki; zwykły obiekt \texttt{std::string}
nie jest w C++20 strukturalnym argumentem szablonu. Rozmiar i
zawartość napisu wpływają tutaj na tożsamość specjalizacji szablonu.

\subsection{Nazwy zależne: \texttt{typename} i \texttt{template}}

W definicji szablonu kompilator nie zawsze wie, czy nazwa zależna od
parametru jest typem, czy wartością. W wielu kontekstach musimy to
jednoznacznie zaznaczyć:

\begin{lstlisting}
#include <vector>

template<class Container>
auto first_copy(const Container& c) {
    typename Container::value_type x = c.front();
    return x;
}

template<class T>
struct Wrapper {
    template<class U>
    static constexpr bool same = false;
};

template<class T>
constexpr bool use_member_template() {
    return Wrapper<T>::template same<int>;
}
\end{lstlisting}

Słowo \texttt{typename} wskazuje, że zależna nazwa członkowska
jest typem. Zapis \texttt{::template same<int>} informuje,
że po operatorze \texttt{::} występuje szablon składowy, a znaki
\texttt{<} i \texttt{>} są delimitatorami argumentów szablonu.
Dokładne reguły zależą od kontekstu; obu słów nie należy dopisywać
mechanicznie wszędzie.

\section{Metafunkcje i \texttt{std::integral\_constant}}
\label{sec:cpp9-integral-constant}

Poniższa silnia obliczana przez szablon pokazuje starszy styl
metaprogramowania. Gdy wynikiem ma być zwykła liczba, funkcja
\texttt{constexpr} jest dziś zazwyczaj prostsza. Szablony klas
pozostają potrzebne zwłaszcza wtedy, gdy wynikiem obliczenia ma być
\emph{typ}.

Klasyczne metaprogramowanie szablonowe opierało się na klasach, które
\emph{przenosiły wyniki obliczenia w swoich składowych statycznych}.
Dla przykładu silnię można wyznaczyć przez rekurencyjne instancjowanie
szablonów:

\begin{lstlisting}
template<unsigned N>
struct Factorial {
    static constexpr unsigned value =
        N * Factorial<N - 1>::value;
};

template<>
struct Factorial<0> {
    static constexpr unsigned value = 1;
};

static_assert(Factorial<5>::value == 120);
\end{lstlisting}

Specjalizacja dla zera pełni rolę warunku stopu. Bez niej kompilator
próbowałby tworzyć kolejne instancje, aż osiągnąłby limit głębokości.
Obliczenie ma charakter rekursji \emph{w strukturze typów}, a nie
rekursji wywołań funkcji wykonywanych w czasie działania programu.

Standardowym fundamentem takich metafunkcji jest
\texttt{std::integral\_constant<T,v>}:

\begin{lstlisting}
#include <type_traits>

using TrueType = std::integral_constant<bool, true>;
static_assert(TrueType::value);
static_assert(std::bool_constant<true>::value);
static_assert(std::true_type::value);
static_assert(!std::false_type::value);

template<unsigned N>
struct PowerOfTwo :
    std::integral_constant<unsigned, 2 * PowerOfTwo<N-1>::value> {};

template<>
struct PowerOfTwo<0> : std::integral_constant<unsigned, 1> {};

static_assert(PowerOfTwo<5>::value == 32);
\end{lstlisting}

Predykaty takie jak \texttt{std::is\_pointer} udostępniają wartość
\texttt{::value}. Od C++17 można używać krótszej wersji
z przyrostkiem \texttt{\_v}, na przykład:
\begin{center}
\texttt{std::is\_pointer\_v<T>}.
\end{center}

\subsection{Metafunkcja szablonowa czy funkcja \texttt{constexpr}?}

\begin{lstlisting}
constexpr unsigned factorial(unsigned n) {
    unsigned result = 1;
    for (unsigned i = 2; i <= n; ++i)
        result *= i;
    return result;
}

static_assert(factorial(5) == 120);
unsigned x = 6;
unsigned y = factorial(x); // allowed at run time
\end{lstlisting}

Jeśli wynik jest zwykłą wartością, funkcja \texttt{constexpr} jest
najczęściej prostsza niż rekurencyjna metafunkcja. Klasyczne metafunkcje
pozostają użyteczne, gdy \emph{wynikiem jest typ}, gdy korzystamy ze
specjalizacji lub interoperujemy ze starszym API.

\section{Specjalizacje: pełne i częściowe}
\label{sec:cpp9-specialization}

Szablon podstawowy definiuje przypadek ogólny. \emph{Specjalizacja pełna}
dotyczy konkretnego zestawu argumentów, a \emph{specjalizacja częściowa}
obsługuje pewną ich rodzinę:

\begin{lstlisting}
#include <type_traits>

template<class T>
struct is_raw_pointer : std::false_type {};

template<class T>
struct is_raw_pointer<T*> : std::true_type {};

static_assert(!is_raw_pointer<int>::value);
static_assert(is_raw_pointer<double*>::value);
static_assert(!is_raw_pointer<int* const>::value);
\end{lstlisting}

Ostatni wynik nie jest pomyłką: \texttt{int* const} jest
\emph{stałym wskaźnikiem}, nie typem postaci \texttt{T*}.
Jeśli chcemy ignorować kwalifikatory zewnętrzne, możemy najpierw
użyć \texttt{std::remove\_cv\_t<T>}.

\begin{lstlisting}
template<class T>
inline constexpr bool is_any_raw_pointer_v =
    is_raw_pointer<std::remove_cv_t<T>>::value;

static_assert(is_any_raw_pointer_v<int* const>);
\end{lstlisting}

\paragraph{Ważne ograniczenie.}
Dla \textbf{szablonów funkcji nie ma specjalizacji częściowej}.
W takich przypadkach wykorzystujemy przeciążanie, koncepcje,
\texttt{if constexpr} lub specjalizację pomocniczego szablonu klasy.
Pełna specjalizacja funkcji istnieje, ale jej interakcja z przeciążeniami
bywa zaskakująca; w nowym kodzie zwykle preferujemy przeciążenia.

\section{\texttt{type\_traits}: pytania o typy}
\label{sec:cpp9-traits}

Biblioteka \texttt{<type\_traits>} jest zbiorem standardowych metafunkcji.
Dzielimy je na \emph{predykaty} (wynik logiczny) i \emph{transformacje}
(wynik będący typem). Ich składnia opiera się na ustalonych konwencjach:
\begin{itemize}
\item \texttt{Trait<T>::value} lub \texttt{Trait\_v<T>} dla predykatów;
\item \texttt{typename Trait<T>::type} lub \texttt{Trait\_t<T>}
      dla transformacji.
\end{itemize}

\begin{lstlisting}
#include <type_traits>
#include <string>

static_assert(std::is_integral_v<int>);
static_assert(std::is_floating_point_v<double>);
static_assert(std::is_same_v<int, int>);
static_assert(std::is_pointer_v<const char*>);
static_assert(std::is_constructible_v<std::string, const char*>);
static_assert(std::is_move_constructible_v<std::string>);

using A = std::remove_const_t<const int>; // int
using B = std::remove_reference_t<int&>; // int
using C = std::remove_cvref_t<const int&>; // int (C++20)
using D = std::add_pointer_t<int>;        // int*
static_assert(std::is_same_v<A, B>);
static_assert(std::is_same_v<C, int>);
static_assert(std::is_same_v<D, int*>);
\end{lstlisting}

\begin{center}
\small
\begin{tabular}{|>{\raggedright\arraybackslash}p{5.5cm}|>{\raggedright\arraybackslash}p{7.0cm}|}
\hline
\textbf{Cecha} & \textbf{Co sprawdza lub wyznacza?} \\
\hline
\texttt{is\_same\_v<A,B>} & Identyczność typów (wraz z istotnymi kwalifikatorami). \\
\hline
\texttt{is\_convertible\_v<A,B>} & Dostępność konwersji niejawnej \texttt{A} do \texttt{B}. \\
\hline
\texttt{is\_constructible\_v<T,Args...>} & Możliwość skonstruowania \texttt{T} z argumentów zadanych typów. \\
\hline
\texttt{is\_nothrow\_move\_}\newline\texttt{constructible\_v<T>} & Deklarowana bezwyjątkowość konstrukcji przenoszącej. \\
\hline
\texttt{is\_trivially\_copyable\_v<T>} & Własność pozwalająca na określone operacje z reprezentacją bajtową. \\
\hline
\texttt{remove\_cvref\_t<T>} & Usunięcie referencji i kwalifikatorów \texttt{const}/\texttt{volatile} z poziomu wierzchniego. \\
\hline
\texttt{common\_type\_t<A,B>} & Wspólny typ wynikowy dla danych typów (o ile istnieje). \\
\hline
\texttt{invoke\_result\_t<F,Args...>} & Typ wyniku wywołania poprzez \texttt{std::invoke}. \\
\hline
\end{tabular}
\end{center}

\paragraph{\texttt{remove\_cvref} a \texttt{decay}.}
\texttt{std::decay\_t<T>} oprócz usunięcia referencji i zewnętrznych
kwalifikatorów może przekształcić tablicę w wskaźnik oraz typ funkcji
w wskaźnik do funkcji. Nie są to zatem synonimy.

\begin{lstlisting}
using X = std::remove_cvref_t<int(&)[3]>; // int[3]
using Y = std::decay_t<int(&)[3]>;        // int*
static_assert(std::is_same_v<Y, int*>);
\end{lstlisting}

\paragraph{Czego predykat nie gwarantuje?}
\texttt{std::is\_move\_constructible\_v<T>} może być prawdziwe również
wtedy, gdy wywołanie z argumentem rvalue używa konstruktora kopiującego.
Podobnie \texttt{is\_trivially\_copyable} nie oznacza, że każdy odczyt
bajtów, rzutowanie lub serializacja do pliku będzie poprawna bądź
przenośna. Cechy typów opisują wybrane własności języka, nie
pełny kontrakt semantyczny.

\subsection{Składanie predykatów z zachowaniem krótkiego obliczania}

\texttt{std::conjunction}, \texttt{std::disjunction} i
\texttt{std::negation} łączą predykaty, które mają składową
\texttt{::value}. Dwa pierwsze narzędzia zatrzymują odczytywanie \texttt{::value}
z kolejnych predykatów, kiedy wynik jest już znany. Trzeba jednak
poprawnie sformować same argumenty typowe tych metafunkcji.

\begin{lstlisting}
#include <type_traits>

template<class T>
using SmallSignedInteger = std::conjunction<
    std::is_integral<T>,
    std::is_signed<T>,
    std::bool_constant<(sizeof(T) <= sizeof(int))>>;

static_assert(SmallSignedInteger<short>::value);
static_assert(!SmallSignedInteger<unsigned>::value);
\end{lstlisting}

Dla prostych stałych logicznych zapis
\texttt{(cond1 \&\& cond2)} bywa czytelniejszy.
\texttt{conjunction} jest użyteczne przede wszystkim wtedy, gdy
poszczególne warunki są odrębnymi metafunkcjami, a ich instancjowanie
może być kosztowne lub niepoprawne bez wcześniejszego warunku.

\section{Metafunkcje zwracające typy i leniwy wybór}
\label{sec:cpp9-transforms}

Tworząc własną transformację typu, zwyczajowo zapisujemy wynik jako
\texttt{using type = ...} i udostępniamy wygodny alias:

\begin{lstlisting}
template<class T>
struct storage_type {
    using type = std::remove_cvref_t<T>;
};

template<class T>
using storage_type_t = typename storage_type<T>::type;

static_assert(std::is_same_v<storage_type_t<const long&>, long>);
\end{lstlisting}

Ważnym narzędziem jest \texttt{std::conditional\_t<B,T,F>}, wybierające
jeden z dwóch typów w zależności od znanej w kompilacji wartości logicznej:

\begin{lstlisting}
template<class T>
using accumulator_t = std::conditional_t<
    std::is_integral_v<T>, long long, long double>;

static_assert(std::is_same_v<accumulator_t<int>, long long>);
static_assert(std::is_same_v<accumulator_t<float>, long double>);
\end{lstlisting}

\paragraph{Pułapka eager instantiation.}
\texttt{conditional\_t} wybiera typ spośród \emph{już utworzonych
argumentów typowych}. Nie chroni przed błędami powstającymi podczas
formowania któregoś z nich:

\begin{lstlisting}
// Incorrect for int: typename T::value_type is invalid.
// template<class T>
// using bad = std::conditional_t<
//     std::is_integral_v<T>, T, typename T::value_type>;
\end{lstlisting}

Aby uzyskać prawdziwie warunkowe tworzenie typu, można zastosować
specjalizację pomocniczej klasy:

\begin{lstlisting}
template<class T, bool = std::is_integral_v<T>>
struct value_or_self {
    using type = typename T::value_type;
};

template<class T>
struct value_or_self<T, true> {
    using type = T;
};

template<class T>
using value_or_self_t = typename value_or_self<T>::type;

static_assert(std::is_same_v<value_or_self_t<int>, int>);
\end{lstlisting}

Podstawowa definicja zawiera nazwę zależną, która jest sprawdzana przy
instancjowaniu odpowiedniego przypadku; nie jest potrzebna dla \texttt{int}.

\section{\texttt{decltype}, \texttt{declval} i konteksty bez obliczania}

\begin{center}\small
\begin{tabularx}{\textwidth}{@{}p{0.32\textwidth}X@{}}
\toprule
Narzędzie & Za co odpowiada \\
\midrule
\code{decltype(expr)} & Ustala typ wyrażenia bez wykonywania go; szczególna reguła dotyczy nieujętej w nawiasy nazwy. \\
\code{decltype((x))} & Dla lvalue \code{x} zwraca typ referencyjny, np. \texttt{T\&}. \\
\code{std::declval<T>()} & Udostępnia fikcyjne wyrażenie typu \code{T} wyłącznie w kontekście nieobliczanym. \\
\code{std::void_t<...>} & Zamienia poprawne podstawienia typów na \code{void}; przydatne w starszych technikach wykrywania. \\
\code{if constexpr} & W szablonie wybiera gałąź podczas kompilacji i odrzuca niepotrzebną gałąź po instancjacji. \\
\bottomrule
\end{tabularx}
\end{center}
\label{sec:cpp9-decltype}

\texttt{decltype(expr)} pozwala odczytać typ wyrażenia bez wykonywania
samego wyrażenia. Subtelność stanowi różnica między nieujętym w nawiasy
identyfikatorem a ogólnym wyrażeniem:

\begin{lstlisting}
int x = 5;
static_assert(std::is_same_v<decltype(x), int>);
static_assert(std::is_same_v<decltype((x)), int&>);
\end{lstlisting}

Wyrażenie \texttt{(x)} jest \textit{lvalue}, więc \texttt{decltype((x))}
zwraca \texttt{int\&}. Reguła ta ma bezpośredni związek z kategoriami
wartości (rozdział~6).

\texttt{std::declval<T>()} pozwala użyć argumentu typu \texttt{T}
w \emph{kontekście niewartościowanym}, nawet jeśli nie możemy utworzyć
jego obiektu. Nie wolno rzeczywiście wykonywać wywołania
\texttt{declval}.

\begin{lstlisting}
#include <utility>
#include <type_traits>

template<class T>
using dereference_t = decltype(*std::declval<T&>());

static_assert(std::is_same_v<dereference_t<int*>, int&>);

struct NonDefault {
    NonDefault() = delete;
    int method() const;
};

using Result = decltype(std::declval<const NonDefault&>().method());
static_assert(std::is_same_v<Result, int>);
\end{lstlisting}

Typ wyniku wywołania można również wyznaczyć za pomocą gotowej metafunkcji:
\begin{center}
\texttt{std::invoke\_result\_t<F, Args...>}.
\end{center}

\section{SFINAE: błędne podstawienie to nie zawsze błąd kompilacji}
\label{sec:cpp9-sfinae}

Ta część pomaga czytać biblioteki napisane przed C++20.
W nowym kodzie wymagania wobec typów zwykle prościej wyrazić
przez \texttt{requires}, opisane w rozdziale~13.

\textbf{SFINAE} (\textit{Substitution Failure Is Not An Error}) mówi,
że niektóre błędy powstałe podczas podstawiania argumentów do szablonu
\emph{nie kończą kompilacji}: dana specjalizacja zostaje wyłączona
z zestawu kandydatów lub odpowiednia specjalizacja częściowa nie pasuje.
Dotyczy to określonych kontekstów bezpośredniego podstawiania;
\textbf{nie} jest mechanizmem przechwytującym dowolne błędy w ciele funkcji
ani skutki uboczne instancjowania innych typów.

\subsection{\texttt{std::enable\_if\_t}}

\begin{lstlisting}
#include <type_traits>

template<class T>
constexpr std::enable_if_t<std::is_integral_v<T>, T>
doubled(T x) {
    return x + x;
}

static_assert(doubled(3) == 6);
// doubled(2.5); // no matching function
\end{lstlisting}

Dla typu \texttt{double} nie istnieje \texttt{type} w
\texttt{std::enable\_if<false,double>}; funkcja zostaje wyłączona
z odpowiedniego zbioru kandydatów. Do konstruktorów, operatorów i
różnych wariantów przeciążeń istnieją odmienne miejsca umieszczania
\texttt{enable\_if}; łatwo przy tym uzyskać przypadkowe powtórne
deklaracje lub nieczytelne błędy.

\subsection{\texttt{void\_t} — detection idiom}

Od C++17 wygodnie wykrywać istnienie nazw i poprawność wyrażeń przy
użyciu \texttt{std::void\_t}. Jeżeli wszystkie podane argumenty
są prawidłowymi typami, rezultatem jest po prostu \texttt{void}:

\begin{lstlisting}
#include <type_traits>
#include <vector>

template<class, class = void>
struct has_value_type : std::false_type {};

template<class T>
struct has_value_type<T, std::void_t<typename T::value_type>>
    : std::true_type {};

static_assert(has_value_type<std::vector<int>>::value);
static_assert(!has_value_type<int>::value);
\end{lstlisting}

W podstawieniu \texttt{T=int} wyrażenie \texttt{int::value\_type}
nie ma sensu; specjalizacja częściowa nie jest więc wybierana.
Analogicznie można badać poprawność wyrażenia z \texttt{decltype}:

\begin{lstlisting}
#include <type_traits>
#include <utility>

template<class, class = void>
struct has_preincrement : std::false_type {};

template<class T>
struct has_preincrement<T,
    std::void_t<decltype(++std::declval<T&>())>>
    : std::true_type {};

static_assert(has_preincrement<int>::value);
\end{lstlisting}

\subsection{SFINAE kontra \texttt{requires} (C++20)}

Konstrukcję wykrywającą operację możemy wyrazić bezpośrednio za pomocą
koncepcji z rozdziału~13:

\begin{lstlisting}
#include <concepts>

template<class T>
concept Preincrementable = requires(T& x) {
    { ++x } -> std::same_as<T&>;
};

template<Preincrementable T>
void advance_one(T& x) { ++x; }
\end{lstlisting}

Metafunkcja \texttt{has\_preincrement<T>} sprawdza jedynie poprawność
operacji \texttt{++x}. Z kolei koncepcja \texttt{Preincrementable}
dodatkowo wymaga, aby jej \emph{typ wyniku} odpowiadał dokładnie \texttt{T\&}.
Różnica ta jest zamierzona: ograniczenia należy definiować zgodnie z
kontraktem, a nie tylko według dostępności pewnego wyrażenia.

\section{\texttt{if constexpr}: wybór implementacji na podstawie typu}
\label{sec:cpp9-ifconstexpr}

Od C++17 często zamiast rekurencyjnych specjalizacji wybieramy gałąź
wewnątrz pojedynczego szablonu. Niebrana gałąź zależna od argumentu
szablonu nie jest instancjowana dla danego zestawu argumentów:

\begin{lstlisting}
#include <type_traits>

template<class T>
constexpr auto number_or_value(const T& x) {
    if constexpr (std::is_arithmetic_v<T>)
        return x;
    else
        return x.value();
}

struct Box {
    constexpr int value() const { return 7; }
};
static_assert(number_or_value(3) == 3);
static_assert(number_or_value(Box{}) == 7);
\end{lstlisting}

\paragraph{Dlaczego nie zwykłe \texttt{if}?}
W zwykłym \texttt{if} kompilator musi przeanalizować obie gałęzie
pod kątem poprawności dla instancjowanego kodu. Dla \texttt{int}
wyrażenie \texttt{x.value()} nie jest poprawne, a sam warunek
\texttt{std::is\_arithmetic\_v<int>} tego nie zmieni.

\paragraph{Ważna subtelność.}
\texttt{if constexpr(false)} nie jest zamiennikiem
\texttt{\#if 0}: poza kontekstem szablonowym odrzucona gałąź nadal
podlega standardowym kontrolom poprawności. Ponadto mechanizm nie
usprawiedliwia konstrukcji zawsze niepoprawnych niezależnie od
argumentów szablonu.

\section{Pakiety typów, sekwencje indeksów i fold expressions}
\label{sec:cpp9-type-lists}

Metaprogramowanie dobrze współpracuje z pakietami parametrów
(rozdział~11). Zamiast tworzyć klasyczne rekurencyjne listy typów,
często możemy używać rozwinięcia pakietu i wyrażenia zwijającego.

\begin{lstlisting}
#include <type_traits>

template<class... Ts>
struct type_list {};

template<class T, class... Ts>
inline constexpr bool contains_v =
    (std::is_same_v<T, Ts> || ...); // empty pack -> false

static_assert(contains_v<int, char, int, double>);
static_assert(!contains_v<long, char, int, double>);
static_assert(!contains_v<int>);
\end{lstlisting}

Zamiast tworzyć specjalizacje dla każdego elementu pakietu, korzystamy
z jednego \texttt{fold expression}. Gdy potrzebna jest reprezentacja
indeksów \texttt{0,1,...,N-1}, biblioteka udostępnia
\texttt{std::make\_index\_sequence<N>}:

\begin{lstlisting}
#include <tuple>
#include <utility>
#include <cstddef>

template<class Tuple, std::size_t... I>
constexpr auto reversed_impl(const Tuple& t,
                             std::index_sequence<I...>) {
    return std::make_tuple(std::get<sizeof...(I)-1-I>(t)...);
}

template<class... Ts>
constexpr auto reversed(const std::tuple<Ts...>& t) {
    return reversed_impl(t, std::index_sequence_for<Ts...>{});
}

static_assert(reversed(std::tuple{1, 2, 3}) ==
              std::tuple{3, 2, 1});
\end{lstlisting}

\paragraph{Czy zawsze potrzebujemy indeksów?}
Nie. W przypadku prostego wykonania funkcji na elementach krotki
w C++17 często wystarczy \texttt{std::apply}; sekwencja indeksów
jest potrzebna zwłaszcza wtedy, gdy chcemy \emph{przekształcić pozycje}
lub mieć indywidualny dostęp do każdego indeksu na etapie kompilacji.

\section{\texttt{constexpr}, \texttt{consteval} i obliczenia na danych}
\label{sec:cpp9-constant-eval}

Szablony nie są jedynym sposobem wykonania pracy w trakcie kompilacji.
Od C++11 rozwijano możliwości \texttt{constexpr}, a w C++20 mamy
również \texttt{consteval}. Różnica jest istotna:
\begin{itemize}
\item funkcja \texttt{constexpr} \emph{może} uczestniczyć w obliczeniach
stałych, a w dozwolonych wywołaniach działa też podczas wykonania;
\item wywołanie funkcji \texttt{consteval} musi spełniać reguły
właściwe obliczaniu w czasie kompilacji (\textit{immediate function}).
\end{itemize}

\begin{lstlisting}
constexpr int square(int x) { return x * x; }
consteval int cube_ct(int x) { return x * x * x; }

static_assert(square(4) == 16);
static_assert(cube_ct(3) == 27);

int n = 4;
int a = square(n); // runtime evaluation possible
// int b = cube_ct(n); // error: n is not constant
\end{lstlisting}

Przykładem praktycznym jest haszowanie \texttt{std::string\_view}
w stałym wyrażeniu. Poniższy wariant FNV-1a nadaje się do zbudowania
stałej identyfikującej np. nazwę polecenia. Nie jest to hash
kryptograficzny i nie wyklucza kolizji.

\begin{lstlisting}
#include <cstdint>
#include <string_view>

consteval std::uint64_t fnv1a(std::string_view text) {
    std::uint64_t h = 14695981039346656037ull;
    for (unsigned char c : text) {
        h ^= c;
        h *= 1099511628211ull;
    }
    return h;
}

static_assert(fnv1a("abc") == 0xe71fa2190541574bull);
\end{lstlisting}

\paragraph{\texttt{constinit} i \texttt{std::is\_constant\_evaluated}.}
\texttt{constinit} zapewnia statyczną inicjalizację odpowiedniej
zmiennej, lecz sam nie czyni jej stałą. Funkcja
\texttt{std::is\_constant\_evaluated()} (C++20) pozwala w
pewnych przypadkach rozpoznać, czy bieżąca ewaluacja ma charakter
stały; nadużywanie jej może jednak zaciemniać semantykę funkcji.

\section{CRTP: polimorfizm statyczny}
\label{sec:cpp9-crtp}

\textbf{CRTP} (\textit{Curiously Recurring Template Pattern}) polega na
tym, że klasa dziedziczy po specjalizacji szablonu, której argumentem
jest ona sama:

\begin{lstlisting}
template<class Derived>
struct Printable {
    void print() const {
        static_cast<const Derived&>(*this).print_impl();
    }
};

struct Number : Printable<Number> {
    int value = 0;
    void print_impl() const;
};
\end{lstlisting}

Mechanizm \texttt{static\_cast} jest tutaj poprawny pod założeniem,
że obiekt bazowy faktycznie stanowi podobiekt \texttt{Derived}.
CRTP nie zapewnia dynamicznego sprawdzania typu ani wirtualnego
wywołania — poprawność wynika z konstrukcji hierarchii. Wywołanie
\texttt{print\_impl()} może być statycznie związane i często
łatwiej je wbudować w kod wynikowy (\textit{inline}).

Bardziej praktycznym przykładem jest domieszka (\textit{mixin})
dostarczająca wspólną operację, lecz działająca na danych
konkretnej klasy pochodnej:

\begin{lstlisting}
template<class Derived>
struct TranslationMixin {
    constexpr void translate(int dx, int dy) {
        auto& self = static_cast<Derived&>(*this);
        self.x += dx;
        self.y += dy;
    }
};

struct Point : TranslationMixin<Point> {
    int x{}, y{};
    constexpr Point(int a, int b) : x(a), y(b) {}
};

constexpr bool test_translation() {
    Point p{1, 2};
    p.translate(3, 4);
    return p.x == 4 && p.y == 6;
}
static_assert(test_translation());
\end{lstlisting}

Baza \texttt{TranslationMixin<Point>} jest instancjowana z typem
\texttt{Point}, chociaż sam typ \texttt{Point} nie jest jeszcze
kompletny w momencie deklarowania dziedziczenia. Odwołania do składowych
\texttt{x} i \texttt{y} zależą jednak od \texttt{Derived} i są
sprawdzane dopiero podczas instancjowania odpowiedniej funkcji.
W rzeczywistym API warto opisać wymagane składowe koncepcją lub
przenieść operację do jawnie wskazanego interfejsu.

\paragraph{CRTP a wirtualne funkcje.}
Klasyczny polimorfizm dynamiczny pozwala wybierać implementację
na podstawie rzeczywistego typu obiektu przez interfejs klasy bazowej,
także gdy typ jest nieznany podczas kompilacji kodu wywołującego.
CRTP wymaga znajomości konkretnego typu podczas instancjowania.
Nie jest ogólnym zamiennikiem klas abstrakcyjnych. W C++20 prostsze
zastosowania domieszek i interfejsów statycznych można często oprzeć
na \texttt{concepts} i funkcjach wolnostojących.

\section{Policy-based design: wybór zachowania przez typ}
\label{sec:cpp9-policy}

Szablon może przyjąć \emph{politykę}, której metody definiują wariant
algorytmu. Umożliwia to wybór zachowania bez wirtualnych wywołań,
a kompilator może całkowicie usunąć niepotrzebne operacje:

\begin{lstlisting}
#include <limits>
#include <stdexcept>

struct Checked {
    static int divide(int a, int b) {
        if (b == 0) throw std::invalid_argument("zero");
        if (a == std::numeric_limits<int>::min() && b == -1)
            throw std::overflow_error("result outside int");
        return a / b;
    }
};

struct Fast {
    static int divide(int a, int b) { return a / b; }
};

template<class Policy>
struct Calculator {
    int divide(int a, int b) const {
        return Policy::divide(a, b);
    }
};

// Calculator<Checked> checks both invalid integer cases;
// Calculator<Fast> requires a nonzero divisor and no overflow.
\end{lstlisting}

Polityka \texttt{Fast} nie zabezpiecza dzielenia przez zero ani
przepełnienia \texttt{INT\_MIN / -1}; oba przypadki mają
niezdefiniowane zachowanie dla \texttt{int}. \texttt{Checked}
odrzuca je przed dzieleniem. Dobór polityki jest częścią kontraktu API.

W realnym kodzie szablon polityki można ograniczyć koncepcją, a jeśli
liczba wariantów jest niewielka, zwykłe \texttt{if constexpr} może być
łatwiejsze w utrzymaniu niż rozbudowana hierarchia polityk.

\section{Kompletny przykład: dobór typu akumulatora}
\label{sec:cpp9-worked}

Załóżmy, że chcemy obliczać sumę liczb, lecz wybierać typ akumulatora
w zależności od argumentu: dla typów całkowitych \texttt{long long},
a dla zmiennoprzecinkowych \texttt{long double}. Rozwiązanie łączy
\texttt{concepts}, \texttt{type\_traits} i \texttt{conditional\_t}:

\begin{lstlisting}
#include <concepts>
#include <type_traits>
#include <vector>

template<class T>
concept Numeric = std::integral<T> || std::floating_point<T>;

template<Numeric T>
using sum_type_t = std::conditional_t<
    std::integral<T>, long long, long double>;

template<Numeric T>
sum_type_t<T> accumulate_values(const std::vector<T>& xs) {
    sum_type_t<T> result{};
    for (const T& x : xs) {
        result += static_cast<sum_type_t<T>>(x);
    }
    return result;
}

static_assert(std::is_same_v<sum_type_t<int>, long long>);
static_assert(std::is_same_v<sum_type_t<double>, long double>);
\end{lstlisting}

\paragraph{Ograniczenia.}
Przykład wyjaśnia selekcję typu, lecz \emph{nie daje gwarancji
braku przepełnienia}. Nawet \texttt{long long} ma ograniczony zakres,
a suma zmiennoprzecinkowa może akumulować błędy zaokrągleń.
Najpierw należy ustalić wymagania numeryczne, a dopiero potem
projektować odpowiednią metafunkcję.

\section{Czas kompilacji, diagnostyka i praktyczne ograniczenia}

Metaprogramowanie przenosi koszt części pracy z czasu wykonania na
czas kompilacji, ale samo w sobie nie oznacza wydajniejszego programu.
Każda odmienna specjalizacja szablonu może wymagać osobnego
instancjowania i wygenerowania kodu. Powstają wtedy m.in.:
\begin{itemize}
\item dłuższe kompilacje i większe zużycie pamięci kompilatora;
\item \emph{code bloat}, czyli powiększenie plików wynikowych przez wiele
podobnych instancji;
\item skomplikowane komunikaty błędów przy zagnieżdżonych metafunkcjach;
\item problemy z zależnościami między plikami nagłówkowymi;
\item ograniczenia głębokości instancjowania i kosztu ewaluacji stałej.
\end{itemize}

\paragraph{Pliki nagłówkowe i jawne instancjowanie.}
Definicje szablonów zazwyczaj umieszczamy w nagłówkach, ponieważ
kompilator musi je znać przy instancjowaniu z konkretnymi argumentami.
Możliwe jest również \emph{jawne instancjowanie} dla ustalonych typów,
co pozwala kontrolować zestaw generowanych specjalizacji:

\begin{lstlisting}
// In a header:
template<class T> T triple(T x);
extern template int triple<int>(int);

// In one .cpp file (definition must be visible here):
template<class T> T triple(T x) { return x + x + x; }
template int triple<int>(int);
\end{lstlisting}

Ten przykład przedstawia ideę rozdzielenia deklaracji i jawnego
instancjowania; obie oznaczone części należy umieścić w odpowiednich
plikach, a dla wywołań innych typów zapewnić dostęp do definicji.
Warto pamiętać o regule ODR z rozdziału~4.

\section{Jaką technikę wybrać?}

\begin{center}
\small
\begin{tabular}{|p{5.7cm}|p{6.8cm}|}
\hline
\textbf{Potrzeba} & \textbf{Najczęściej odpowiedni mechanizm} \\
\hline
Sprawdzenie wymaganych operacji w nowym API & \texttt{concepts}, \texttt{requires} (C++20) \\
\hline
Obliczenie wartości dla stałych danych & Funkcja \texttt{constexpr} \\
\hline
Wymuszenie ewaluacji w czasie kompilacji & \texttt{consteval} (C++20) \\
\hline
Warunkowy kod zależny od typu & \texttt{if constexpr} \\
\hline
Ustalenie własności typu & \texttt{std::is\_...\_v}, inne \texttt{type\_traits} \\
\hline
Zmiana lub wybór typu wynikowego & \texttt{...\_t}, specjalizacja pomocniczej klasy \\
\hline
Wsparcie biblioteki sprzed C++20 & SFINAE, \texttt{enable\_if}, \texttt{void\_t} \\
\hline
Wieloargumentowe cechy typów & Pakiety typów, \texttt{fold expressions} \\
\hline
Dostęp do pozycji elementów krotki & \texttt{index\_sequence}, \texttt{std::apply} \\
\hline
Wielokrotne użycie zachowania klasowego & CRTP / mixin, rozważ też koncepcje \\
\hline
Wybór wariantu strategii w typie & Szablon polityki (policy-based design) \\
\hline
\end{tabular}
\end{center}

\section{Typowe błędy i zasady projektowe}

\begin{enumerate}
\item \textbf{Mylenie typu z wartością.} \texttt{typename T::value\_type}
      to typ, a \texttt{Trait<T>::value} to wartość; nie są wymienne.
\item \textbf{Brak warunku stopu.} Rekurencyjny szablon klasy może
      prowadzić do nieograniczonego instancjowania i błędów głębokości.
\item \textbf{Założenie, że SFINAE przechwytuje wszystkie błędy.}
      Działa tylko w określonych kontekstach podstawiania.
\item \textbf{Nadużywanie \texttt{conditional\_t}.} Oba argumenty
      typowe trzeba poprawnie sformować, zanim zostanie dokonany wybór.
\item \textbf{Próba wywołania \texttt{declval}.} Można go używać wyłącznie
      w odpowiednich kontekstach niewartościowanych.
\item \textbf{Mylenie \texttt{decltype(x)} z \texttt{decltype((x))}.}
      Dodatkowe nawiasy mogą zmienić wynik na typ referencyjny.
\item \textbf{Mylenie \texttt{constexpr} z obowiązkiem obliczenia
      podczas kompilacji.} Nie każde wywołanie funkcji \texttt{constexpr}
      jest ewaluowane w czasie kompilacji.
\item \textbf{Przenoszenie semantyki do traits.} Typ spełniający
      syntaktyczne wymaganie nie musi spełniać własności biznesowej,
      np. braku utraty precyzji czy bezpiecznej serializacji.
\item \textbf{Niepotrzebny CRTP.} Gdy wystarczy wolna funkcja
      z koncepcją, hierarchia klas często niczego nie upraszcza.
\item \textbf{Przedwczesne metaprogramowanie.} Najprostsze czytelne
      rozwiązanie jest zazwyczaj lepsze od dodatkowej warstwy abstrakcji.
\end{enumerate}

\paragraph{Pytania kontrolne.}
\begin{enumerate}
\item Czym różni się instancjowanie szablonu od wykonania funkcji?
\item Dlaczego \texttt{Factorial<0>} wymaga osobnej specjalizacji?
\item Kiedy \texttt{std::decay\_t<T>} i \texttt{std::remove\_cvref\_t<T>}
      dają różne wyniki?
\item Dlaczego \texttt{std::conditional\_t} nie rozwiązuje wszystkich
      problemów warunkowego tworzenia typów?
\item Co dokładnie sprawdza \texttt{void\_t<decltype(++declval<T\&>())>}?
\item Kiedy niepowodzenie podstawienia jest SFINAE, a kiedy zwykłym
      błędem kompilacji?
\item Dlaczego \texttt{if constexpr} nie zastępuje preprocesora?
\item Czym różnią się \texttt{constexpr} i \texttt{consteval}?
\item W jakich sytuacjach CRTP nie zastąpi funkcji wirtualnych?
\item Kiedy lepiej zastosować \texttt{concepts} niż \texttt{enable\_if}?
\end{enumerate}

\paragraph{Dalsza lektura.}
Warto sięgnąć do dokumentacji \emph{cppreference} dotyczącej
\emph{Templates}, \emph{Type traits} i \emph{SFINAE}.
Pomocne są również hasła \texttt{std::void\_t}, \texttt{std::enable\_if},
\texttt{if constexpr}, \texttt{consteval} oraz
\texttt{std::integer\_sequence}. Do diagnozowania złożonych
błędów szablonów przydają się opcje \texttt{-Wall -Wextra -Wpedantic},
\texttt{-ftemplate-backtrace-limit=0} (GCC) oraz ograniczanie zasięgu
instancjowanego kodu do minimalnego przykładu.


\chapter[Metaprogramowanie w C++]{Zaawansowane metaprogramowanie: szablony, cechy typów i obliczenia podczas kompilacji}
\label{chap:meta-cpp}

\section{Czym jest metaprogramowanie?}

\emph{Metaprogramowanie} oznacza pisanie kodu, który na etapie kompilacji
wyznacza typy, własności typów, wartości lub strukturę później
kompilowanego programu. Szablon nie jest gotową funkcją czy klasą:
stanowi przepis, z którego kompilator tworzy konkretne wersje
dla podanych typów. Nazywamy to \emph{instancjowaniem szablonu}.
Na przykład wywołanie
\texttt{twice(21)} w programie poniżej każe kompilatorowi przygotować
wersję funkcji dla \texttt{int}; program podczas działania wywołuje
już tę wersję. W C++20 do dyspozycji mamy kilka narzędzi:
\begin{itemize}
\item \textbf{szablony klas i funkcji} — generowanie kodu dla różnych typów;
\item \textbf{\texttt{type\_traits}} — obliczanie własności i przekształcanie typów;
\item \textbf{specjalizacje i SFINAE} — wybór poprawnych wersji szablonów;
\item \textbf{\texttt{if constexpr}, \texttt{constexpr}, \texttt{consteval}} —
      wykonywanie obliczeń i selekcja kodu w czasie kompilacji;
\item \textbf{\texttt{concepts}} — czytelne formułowanie wymagań wobec typów;
\item \textbf{CRTP} — wykorzystanie dziedziczenia do polimorfizmu statycznego.
\end{itemize}

Metaprogramowanie ma praktyczne cele: eliminowanie powtarzalnego kodu,
uzyskiwanie lepszych błędów kompilacji, wybór wydajnej implementacji oraz
zapewnienie poprawności określonych operacji \emph{przed uruchomieniem} programu.
Nie oznacza automatycznie zwiększenia wydajności — kosztem mogą być dłuższa
kompilacja, powiększenie kodu maszynowego i trudniejsze komunikaty diagnostyczne.

\begin{lstlisting}
#include <type_traits>

template<class T>
T twice(T x) {
    static_assert(std::is_arithmetic_v<T>);
    return x + x;
}

int a = twice(21);       // T = int
// twice("text");       // static_assert fails
\end{lstlisting}

\paragraph{Związki z innymi rozdziałami.}
Wykorzystamy szablony wariadyczne (rozdział~11),
przenoszenie i referencje (rozdział~6) oraz lambdy (rozdział~9).
Koncepcje szerzej poznamy w rozdziale~13; tutaj pokazujemy
przede wszystkim mechanizmy, na których się opierają.
Wszystkie oznaczone jako poprawne przykłady odnoszą się do \textbf{C++20}.

\section{Instancjowanie szablonów, parametry i zależne nazwy}
\label{sec:cpp9-instantiation}

\subsection{Parametry typowe, nietypowe i szablonowe}

Szablon może przyjmować argumenty będące typami, wartościami stałymi
lub innymi szablonami. Parametr wartości (\textit{non-type template parameter},
NTTP) w C++20 może również mieć odpowiedni typ klasowy spełniający
wymagania tzw. typu strukturalnego.

\begin{lstlisting}
#include <array>
#include <cstddef>

template<class T, std::size_t N>
struct Buffer {
    std::array<T, N> data{};
};

Buffer<double, 4> samples; // T=double, N=4
\end{lstlisting}

\texttt{Buffer<double,4>} i \texttt{Buffer<double,8>} są różnymi typami.
Parametr \texttt{N} jest znany podczas kompilacji i może wystąpić tam,
gdzie język wymaga wyrażenia stałego, np. jako rozmiar \texttt{std::array}.

\subsection{Parametr będący obiektem klasy (C++20)}

Od C++20 odpowiednie \emph{typy strukturalne} mogą występować jako
parametry szablonu będące wartościami. Umożliwia to np. identyfikowanie
szablonu przez napis znany podczas kompilacji:

\begin{lstlisting}
#include <cstddef>

template<std::size_t N>
struct fixed_string {
    char data[N]{};
    constexpr fixed_string(const char (&s)[N]) {
        for (std::size_t i = 0; i < N; ++i)
            data[i] = s[i];
    }
};

template<fixed_string Name>
struct NamedCommand {
    static constexpr auto name = Name;
};

static_assert(NamedCommand<"run">::name.data[0] == 'r');
\end{lstlisting}

To przykład zaawansowanej techniki; zwykły obiekt \texttt{std::string}
nie jest w C++20 strukturalnym argumentem szablonu. Rozmiar i
zawartość napisu wpływają tutaj na tożsamość specjalizacji szablonu.

\subsection{Nazwy zależne: \texttt{typename} i \texttt{template}}

W definicji szablonu kompilator nie zawsze wie, czy nazwa zależna od
parametru jest typem, czy wartością. W wielu kontekstach musimy to
jednoznacznie zaznaczyć:

\begin{lstlisting}
#include <vector>

template<class Container>
auto first_copy(const Container& c) {
    typename Container::value_type x = c.front();
    return x;
}

template<class T>
struct Wrapper {
    template<class U>
    static constexpr bool same = false;
};

template<class T>
constexpr bool use_member_template() {
    return Wrapper<T>::template same<int>;
}
\end{lstlisting}

Słowo \texttt{typename} wskazuje, że zależna nazwa członkowska
jest typem. Zapis \texttt{::template same<int>} informuje,
że po operatorze \texttt{::} występuje szablon składowy, a znaki
\texttt{<} i \texttt{>} są delimitatorami argumentów szablonu.
Dokładne reguły zależą od kontekstu; obu słów nie należy dopisywać
mechanicznie wszędzie.

\section{Metafunkcje i \texttt{std::integral\_constant}}
\label{sec:cpp9-integral-constant}

Poniższa silnia obliczana przez szablon pokazuje starszy styl
metaprogramowania. Gdy wynikiem ma być zwykła liczba, funkcja
\texttt{constexpr} jest dziś zazwyczaj prostsza. Szablony klas
pozostają potrzebne zwłaszcza wtedy, gdy wynikiem obliczenia ma być
\emph{typ}.

Klasyczne metaprogramowanie szablonowe opierało się na klasach, które
\emph{przenosiły wyniki obliczenia w swoich składowych statycznych}.
Dla przykładu silnię można wyznaczyć przez rekurencyjne instancjowanie
szablonów:

\begin{lstlisting}
template<unsigned N>
struct Factorial {
    static constexpr unsigned value =
        N * Factorial<N - 1>::value;
};

template<>
struct Factorial<0> {
    static constexpr unsigned value = 1;
};

static_assert(Factorial<5>::value == 120);
\end{lstlisting}

Specjalizacja dla zera pełni rolę warunku stopu. Bez niej kompilator
próbowałby tworzyć kolejne instancje, aż osiągnąłby limit głębokości.
Obliczenie ma charakter rekursji \emph{w strukturze typów}, a nie
rekursji wywołań funkcji wykonywanych w czasie działania programu.

Standardowym fundamentem takich metafunkcji jest
\texttt{std::integral\_constant<T,v>}:

\begin{lstlisting}
#include <type_traits>

using TrueType = std::integral_constant<bool, true>;
static_assert(TrueType::value);
static_assert(std::bool_constant<true>::value);
static_assert(std::true_type::value);
static_assert(!std::false_type::value);

template<unsigned N>
struct PowerOfTwo :
    std::integral_constant<unsigned, 2 * PowerOfTwo<N-1>::value> {};

template<>
struct PowerOfTwo<0> : std::integral_constant<unsigned, 1> {};

static_assert(PowerOfTwo<5>::value == 32);
\end{lstlisting}

Predykaty takie jak \texttt{std::is\_pointer} udostępniają wartość
\texttt{::value}. Od C++17 można używać krótszej wersji
z przyrostkiem \texttt{\_v}, na przykład:
\begin{center}
\texttt{std::is\_pointer\_v<T>}.
\end{center}

\subsection{Metafunkcja szablonowa czy funkcja \texttt{constexpr}?}

\begin{lstlisting}
constexpr unsigned factorial(unsigned n) {
    unsigned result = 1;
    for (unsigned i = 2; i <= n; ++i)
        result *= i;
    return result;
}

static_assert(factorial(5) == 120);
unsigned x = 6;
unsigned y = factorial(x); // allowed at run time
\end{lstlisting}

Jeśli wynik jest zwykłą wartością, funkcja \texttt{constexpr} jest
najczęściej prostsza niż rekurencyjna metafunkcja. Klasyczne metafunkcje
pozostają użyteczne, gdy \emph{wynikiem jest typ}, gdy korzystamy ze
specjalizacji lub interoperujemy ze starszym API.

\section{Specjalizacje: pełne i częściowe}
\label{sec:cpp9-specialization}

Szablon podstawowy definiuje przypadek ogólny. \emph{Specjalizacja pełna}
dotyczy konkretnego zestawu argumentów, a \emph{specjalizacja częściowa}
obsługuje pewną ich rodzinę:

\begin{lstlisting}
#include <type_traits>

template<class T>
struct is_raw_pointer : std::false_type {};

template<class T>
struct is_raw_pointer<T*> : std::true_type {};

static_assert(!is_raw_pointer<int>::value);
static_assert(is_raw_pointer<double*>::value);
static_assert(!is_raw_pointer<int* const>::value);
\end{lstlisting}

Ostatni wynik nie jest pomyłką: \texttt{int* const} jest
\emph{stałym wskaźnikiem}, nie typem postaci \texttt{T*}.
Jeśli chcemy ignorować kwalifikatory zewnętrzne, możemy najpierw
użyć \texttt{std::remove\_cv\_t<T>}.

\begin{lstlisting}
template<class T>
inline constexpr bool is_any_raw_pointer_v =
    is_raw_pointer<std::remove_cv_t<T>>::value;

static_assert(is_any_raw_pointer_v<int* const>);
\end{lstlisting}

\paragraph{Ważne ograniczenie.}
Dla \textbf{szablonów funkcji nie ma specjalizacji częściowej}.
W takich przypadkach wykorzystujemy przeciążanie, koncepcje,
\texttt{if constexpr} lub specjalizację pomocniczego szablonu klasy.
Pełna specjalizacja funkcji istnieje, ale jej interakcja z przeciążeniami
bywa zaskakująca; w nowym kodzie zwykle preferujemy przeciążenia.

\section{\texttt{type\_traits}: pytania o typy}
\label{sec:cpp9-traits}

Biblioteka \texttt{<type\_traits>} jest zbiorem standardowych metafunkcji.
Dzielimy je na \emph{predykaty} (wynik logiczny) i \emph{transformacje}
(wynik będący typem). Ich składnia opiera się na ustalonych konwencjach:
\begin{itemize}
\item \texttt{Trait<T>::value} lub \texttt{Trait\_v<T>} dla predykatów;
\item \texttt{typename Trait<T>::type} lub \texttt{Trait\_t<T>}
      dla transformacji.
\end{itemize}

\begin{lstlisting}
#include <type_traits>
#include <string>

static_assert(std::is_integral_v<int>);
static_assert(std::is_floating_point_v<double>);
static_assert(std::is_same_v<int, int>);
static_assert(std::is_pointer_v<const char*>);
static_assert(std::is_constructible_v<std::string, const char*>);
static_assert(std::is_move_constructible_v<std::string>);

using A = std::remove_const_t<const int>; // int
using B = std::remove_reference_t<int&>; // int
using C = std::remove_cvref_t<const int&>; // int (C++20)
using D = std::add_pointer_t<int>;        // int*
static_assert(std::is_same_v<A, B>);
static_assert(std::is_same_v<C, int>);
static_assert(std::is_same_v<D, int*>);
\end{lstlisting}

\begin{center}
\small
\begin{tabular}{|>{\raggedright\arraybackslash}p{5.5cm}|>{\raggedright\arraybackslash}p{7.0cm}|}
\hline
\textbf{Cecha} & \textbf{Co sprawdza lub wyznacza?} \\
\hline
\texttt{is\_same\_v<A,B>} & Identyczność typów (wraz z istotnymi kwalifikatorami). \\
\hline
\texttt{is\_convertible\_v<A,B>} & Dostępność konwersji niejawnej \texttt{A} do \texttt{B}. \\
\hline
\texttt{is\_constructible\_v<T,Args...>} & Możliwość skonstruowania \texttt{T} z argumentów zadanych typów. \\
\hline
\texttt{is\_nothrow\_move\_}\newline\texttt{constructible\_v<T>} & Deklarowana bezwyjątkowość konstrukcji przenoszącej. \\
\hline
\texttt{is\_trivially\_copyable\_v<T>} & Własność pozwalająca na określone operacje z reprezentacją bajtową. \\
\hline
\texttt{remove\_cvref\_t<T>} & Usunięcie referencji i kwalifikatorów \texttt{const}/\texttt{volatile} z poziomu wierzchniego. \\
\hline
\texttt{common\_type\_t<A,B>} & Wspólny typ wynikowy dla danych typów (o ile istnieje). \\
\hline
\texttt{invoke\_result\_t<F,Args...>} & Typ wyniku wywołania poprzez \texttt{std::invoke}. \\
\hline
\end{tabular}
\end{center}

\paragraph{\texttt{remove\_cvref} a \texttt{decay}.}
\texttt{std::decay\_t<T>} oprócz usunięcia referencji i zewnętrznych
kwalifikatorów może przekształcić tablicę w wskaźnik oraz typ funkcji
w wskaźnik do funkcji. Nie są to zatem synonimy.

\begin{lstlisting}
using X = std::remove_cvref_t<int(&)[3]>; // int[3]
using Y = std::decay_t<int(&)[3]>;        // int*
static_assert(std::is_same_v<Y, int*>);
\end{lstlisting}

\paragraph{Czego predykat nie gwarantuje?}
\texttt{std::is\_move\_constructible\_v<T>} może być prawdziwe również
wtedy, gdy wywołanie z argumentem rvalue używa konstruktora kopiującego.
Podobnie \texttt{is\_trivially\_copyable} nie oznacza, że każdy odczyt
bajtów, rzutowanie lub serializacja do pliku będzie poprawna bądź
przenośna. Cechy typów opisują wybrane własności języka, nie
pełny kontrakt semantyczny.

\subsection{Składanie predykatów z zachowaniem krótkiego obliczania}

\texttt{std::conjunction}, \texttt{std::disjunction} i
\texttt{std::negation} łączą predykaty, które mają składową
\texttt{::value}. Dwa pierwsze narzędzia zatrzymują odczytywanie \texttt{::value}
z kolejnych predykatów, kiedy wynik jest już znany. Trzeba jednak
poprawnie sformować same argumenty typowe tych metafunkcji.

\begin{lstlisting}
#include <type_traits>

template<class T>
using SmallSignedInteger = std::conjunction<
    std::is_integral<T>,
    std::is_signed<T>,
    std::bool_constant<(sizeof(T) <= sizeof(int))>>;

static_assert(SmallSignedInteger<short>::value);
static_assert(!SmallSignedInteger<unsigned>::value);
\end{lstlisting}

Dla prostych stałych logicznych zapis
\texttt{(cond1 \&\& cond2)} bywa czytelniejszy.
\texttt{conjunction} jest użyteczne przede wszystkim wtedy, gdy
poszczególne warunki są odrębnymi metafunkcjami, a ich instancjowanie
może być kosztowne lub niepoprawne bez wcześniejszego warunku.

\section{Metafunkcje zwracające typy i leniwy wybór}
\label{sec:cpp9-transforms}

Tworząc własną transformację typu, zwyczajowo zapisujemy wynik jako
\texttt{using type = ...} i udostępniamy wygodny alias:

\begin{lstlisting}
template<class T>
struct storage_type {
    using type = std::remove_cvref_t<T>;
};

template<class T>
using storage_type_t = typename storage_type<T>::type;

static_assert(std::is_same_v<storage_type_t<const long&>, long>);
\end{lstlisting}

Ważnym narzędziem jest \texttt{std::conditional\_t<B,T,F>}, wybierające
jeden z dwóch typów w zależności od znanej w kompilacji wartości logicznej:

\begin{lstlisting}
template<class T>
using accumulator_t = std::conditional_t<
    std::is_integral_v<T>, long long, long double>;

static_assert(std::is_same_v<accumulator_t<int>, long long>);
static_assert(std::is_same_v<accumulator_t<float>, long double>);
\end{lstlisting}

\paragraph{Pułapka eager instantiation.}
\texttt{conditional\_t} wybiera typ spośród \emph{już utworzonych
argumentów typowych}. Nie chroni przed błędami powstającymi podczas
formowania któregoś z nich:

\begin{lstlisting}
// Incorrect for int: typename T::value_type is invalid.
// template<class T>
// using bad = std::conditional_t<
//     std::is_integral_v<T>, T, typename T::value_type>;
\end{lstlisting}

Aby uzyskać prawdziwie warunkowe tworzenie typu, można zastosować
specjalizację pomocniczej klasy:

\begin{lstlisting}
template<class T, bool = std::is_integral_v<T>>
struct value_or_self {
    using type = typename T::value_type;
};

template<class T>
struct value_or_self<T, true> {
    using type = T;
};

template<class T>
using value_or_self_t = typename value_or_self<T>::type;

static_assert(std::is_same_v<value_or_self_t<int>, int>);
\end{lstlisting}

Podstawowa definicja zawiera nazwę zależną, która jest sprawdzana przy
instancjowaniu odpowiedniego przypadku; nie jest potrzebna dla \texttt{int}.

\section{\texttt{decltype}, \texttt{declval} i konteksty bez obliczania}

\begin{center}\small
\begin{tabularx}{\textwidth}{@{}p{0.32\textwidth}X@{}}
\toprule
Narzędzie & Za co odpowiada \\
\midrule
\code{decltype(expr)} & Ustala typ wyrażenia bez wykonywania go; szczególna reguła dotyczy nieujętej w nawiasy nazwy. \\
\code{decltype((x))} & Dla lvalue \code{x} zwraca typ referencyjny, np. \texttt{T\&}. \\
\code{std::declval<T>()} & Udostępnia fikcyjne wyrażenie typu \code{T} wyłącznie w kontekście nieobliczanym. \\
\code{std::void_t<...>} & Zamienia poprawne podstawienia typów na \code{void}; przydatne w starszych technikach wykrywania. \\
\code{if constexpr} & W szablonie wybiera gałąź podczas kompilacji i odrzuca niepotrzebną gałąź po instancjacji. \\
\bottomrule
\end{tabularx}
\end{center}
\label{sec:cpp9-decltype}

\texttt{decltype(expr)} pozwala odczytać typ wyrażenia bez wykonywania
samego wyrażenia. Subtelność stanowi różnica między nieujętym w nawiasy
identyfikatorem a ogólnym wyrażeniem:

\begin{lstlisting}
int x = 5;
static_assert(std::is_same_v<decltype(x), int>);
static_assert(std::is_same_v<decltype((x)), int&>);
\end{lstlisting}

Wyrażenie \texttt{(x)} jest \textit{lvalue}, więc \texttt{decltype((x))}
zwraca \texttt{int\&}. Reguła ta ma bezpośredni związek z kategoriami
wartości (rozdział~6).

\texttt{std::declval<T>()} pozwala użyć argumentu typu \texttt{T}
w \emph{kontekście niewartościowanym}, nawet jeśli nie możemy utworzyć
jego obiektu. Nie wolno rzeczywiście wykonywać wywołania
\texttt{declval}.

\begin{lstlisting}
#include <utility>
#include <type_traits>

template<class T>
using dereference_t = decltype(*std::declval<T&>());

static_assert(std::is_same_v<dereference_t<int*>, int&>);

struct NonDefault {
    NonDefault() = delete;
    int method() const;
};

using Result = decltype(std::declval<const NonDefault&>().method());
static_assert(std::is_same_v<Result, int>);
\end{lstlisting}

Typ wyniku wywołania można również wyznaczyć za pomocą gotowej metafunkcji:
\begin{center}
\texttt{std::invoke\_result\_t<F, Args...>}.
\end{center}

\section{SFINAE: błędne podstawienie to nie zawsze błąd kompilacji}
\label{sec:cpp9-sfinae}

Ta część pomaga czytać biblioteki napisane przed C++20.
W nowym kodzie wymagania wobec typów zwykle prościej wyrazić
przez \texttt{requires}, opisane w rozdziale~13.

\textbf{SFINAE} (\textit{Substitution Failure Is Not An Error}) mówi,
że niektóre błędy powstałe podczas podstawiania argumentów do szablonu
\emph{nie kończą kompilacji}: dana specjalizacja zostaje wyłączona
z zestawu kandydatów lub odpowiednia specjalizacja częściowa nie pasuje.
Dotyczy to określonych kontekstów bezpośredniego podstawiania;
\textbf{nie} jest mechanizmem przechwytującym dowolne błędy w ciele funkcji
ani skutki uboczne instancjowania innych typów.

\subsection{\texttt{std::enable\_if\_t}}

\begin{lstlisting}
#include <type_traits>

template<class T>
constexpr std::enable_if_t<std::is_integral_v<T>, T>
doubled(T x) {
    return x + x;
}

static_assert(doubled(3) == 6);
// doubled(2.5); // no matching function
\end{lstlisting}

Dla typu \texttt{double} nie istnieje \texttt{type} w
\texttt{std::enable\_if<false,double>}; funkcja zostaje wyłączona
z odpowiedniego zbioru kandydatów. Do konstruktorów, operatorów i
różnych wariantów przeciążeń istnieją odmienne miejsca umieszczania
\texttt{enable\_if}; łatwo przy tym uzyskać przypadkowe powtórne
deklaracje lub nieczytelne błędy.

\subsection{\texttt{void\_t} — detection idiom}

Od C++17 wygodnie wykrywać istnienie nazw i poprawność wyrażeń przy
użyciu \texttt{std::void\_t}. Jeżeli wszystkie podane argumenty
są prawidłowymi typami, rezultatem jest po prostu \texttt{void}:

\begin{lstlisting}
#include <type_traits>
#include <vector>

template<class, class = void>
struct has_value_type : std::false_type {};

template<class T>
struct has_value_type<T, std::void_t<typename T::value_type>>
    : std::true_type {};

static_assert(has_value_type<std::vector<int>>::value);
static_assert(!has_value_type<int>::value);
\end{lstlisting}

W podstawieniu \texttt{T=int} wyrażenie \texttt{int::value\_type}
nie ma sensu; specjalizacja częściowa nie jest więc wybierana.
Analogicznie można badać poprawność wyrażenia z \texttt{decltype}:

\begin{lstlisting}
#include <type_traits>
#include <utility>

template<class, class = void>
struct has_preincrement : std::false_type {};

template<class T>
struct has_preincrement<T,
    std::void_t<decltype(++std::declval<T&>())>>
    : std::true_type {};

static_assert(has_preincrement<int>::value);
\end{lstlisting}

\subsection{SFINAE kontra \texttt{requires} (C++20)}

Konstrukcję wykrywającą operację możemy wyrazić bezpośrednio za pomocą
koncepcji z rozdziału~13:

\begin{lstlisting}
#include <concepts>

template<class T>
concept Preincrementable = requires(T& x) {
    { ++x } -> std::same_as<T&>;
};

template<Preincrementable T>
void advance_one(T& x) { ++x; }
\end{lstlisting}

Metafunkcja \texttt{has\_preincrement<T>} sprawdza jedynie poprawność
operacji \texttt{++x}. Z kolei koncepcja \texttt{Preincrementable}
dodatkowo wymaga, aby jej \emph{typ wyniku} odpowiadał dokładnie \texttt{T\&}.
Różnica ta jest zamierzona: ograniczenia należy definiować zgodnie z
kontraktem, a nie tylko według dostępności pewnego wyrażenia.

\section{\texttt{if constexpr}: wybór implementacji na podstawie typu}
\label{sec:cpp9-ifconstexpr}

Od C++17 często zamiast rekurencyjnych specjalizacji wybieramy gałąź
wewnątrz pojedynczego szablonu. Niebrana gałąź zależna od argumentu
szablonu nie jest instancjowana dla danego zestawu argumentów:

\begin{lstlisting}
#include <type_traits>

template<class T>
constexpr auto number_or_value(const T& x) {
    if constexpr (std::is_arithmetic_v<T>)
        return x;
    else
        return x.value();
}

struct Box {
    constexpr int value() const { return 7; }
};
static_assert(number_or_value(3) == 3);
static_assert(number_or_value(Box{}) == 7);
\end{lstlisting}

\paragraph{Dlaczego nie zwykłe \texttt{if}?}
W zwykłym \texttt{if} kompilator musi przeanalizować obie gałęzie
pod kątem poprawności dla instancjowanego kodu. Dla \texttt{int}
wyrażenie \texttt{x.value()} nie jest poprawne, a sam warunek
\texttt{std::is\_arithmetic\_v<int>} tego nie zmieni.

\paragraph{Ważna subtelność.}
\texttt{if constexpr(false)} nie jest zamiennikiem
\texttt{\#if 0}: poza kontekstem szablonowym odrzucona gałąź nadal
podlega standardowym kontrolom poprawności. Ponadto mechanizm nie
usprawiedliwia konstrukcji zawsze niepoprawnych niezależnie od
argumentów szablonu.

\section{Pakiety typów, sekwencje indeksów i fold expressions}
\label{sec:cpp9-type-lists}

Metaprogramowanie dobrze współpracuje z pakietami parametrów
(rozdział~11). Zamiast tworzyć klasyczne rekurencyjne listy typów,
często możemy używać rozwinięcia pakietu i wyrażenia zwijającego.

\begin{lstlisting}
#include <type_traits>

template<class... Ts>
struct type_list {};

template<class T, class... Ts>
inline constexpr bool contains_v =
    (std::is_same_v<T, Ts> || ...); // empty pack -> false

static_assert(contains_v<int, char, int, double>);
static_assert(!contains_v<long, char, int, double>);
static_assert(!contains_v<int>);
\end{lstlisting}

Zamiast tworzyć specjalizacje dla każdego elementu pakietu, korzystamy
z jednego \texttt{fold expression}. Gdy potrzebna jest reprezentacja
indeksów \texttt{0,1,...,N-1}, biblioteka udostępnia
\texttt{std::make\_index\_sequence<N>}:

\begin{lstlisting}
#include <tuple>
#include <utility>
#include <cstddef>

template<class Tuple, std::size_t... I>
constexpr auto reversed_impl(const Tuple& t,
                             std::index_sequence<I...>) {
    return std::make_tuple(std::get<sizeof...(I)-1-I>(t)...);
}

template<class... Ts>
constexpr auto reversed(const std::tuple<Ts...>& t) {
    return reversed_impl(t, std::index_sequence_for<Ts...>{});
}

static_assert(reversed(std::tuple{1, 2, 3}) ==
              std::tuple{3, 2, 1});
\end{lstlisting}

\paragraph{Czy zawsze potrzebujemy indeksów?}
Nie. W przypadku prostego wykonania funkcji na elementach krotki
w C++17 często wystarczy \texttt{std::apply}; sekwencja indeksów
jest potrzebna zwłaszcza wtedy, gdy chcemy \emph{przekształcić pozycje}
lub mieć indywidualny dostęp do każdego indeksu na etapie kompilacji.

\section{\texttt{constexpr}, \texttt{consteval} i obliczenia na danych}
\label{sec:cpp9-constant-eval}

Szablony nie są jedynym sposobem wykonania pracy w trakcie kompilacji.
Od C++11 rozwijano możliwości \texttt{constexpr}, a w C++20 mamy
również \texttt{consteval}. Różnica jest istotna:
\begin{itemize}
\item funkcja \texttt{constexpr} \emph{może} uczestniczyć w obliczeniach
stałych, a w dozwolonych wywołaniach działa też podczas wykonania;
\item wywołanie funkcji \texttt{consteval} musi spełniać reguły
właściwe obliczaniu w czasie kompilacji (\textit{immediate function}).
\end{itemize}

\begin{lstlisting}
constexpr int square(int x) { return x * x; }
consteval int cube_ct(int x) { return x * x * x; }

static_assert(square(4) == 16);
static_assert(cube_ct(3) == 27);

int n = 4;
int a = square(n); // runtime evaluation possible
// int b = cube_ct(n); // error: n is not constant
\end{lstlisting}

Przykładem praktycznym jest haszowanie \texttt{std::string\_view}
w stałym wyrażeniu. Poniższy wariant FNV-1a nadaje się do zbudowania
stałej identyfikującej np. nazwę polecenia. Nie jest to hash
kryptograficzny i nie wyklucza kolizji.

\begin{lstlisting}
#include <cstdint>
#include <string_view>

consteval std::uint64_t fnv1a(std::string_view text) {
    std::uint64_t h = 14695981039346656037ull;
    for (unsigned char c : text) {
        h ^= c;
        h *= 1099511628211ull;
    }
    return h;
}

static_assert(fnv1a("abc") == 0xe71fa2190541574bull);
\end{lstlisting}

\paragraph{\texttt{constinit} i \texttt{std::is\_constant\_evaluated}.}
\texttt{constinit} zapewnia statyczną inicjalizację odpowiedniej
zmiennej, lecz sam nie czyni jej stałą. Funkcja
\texttt{std::is\_constant\_evaluated()} (C++20) pozwala w
pewnych przypadkach rozpoznać, czy bieżąca ewaluacja ma charakter
stały; nadużywanie jej może jednak zaciemniać semantykę funkcji.

\section{CRTP: polimorfizm statyczny}
\label{sec:cpp9-crtp}

\textbf{CRTP} (\textit{Curiously Recurring Template Pattern}) polega na
tym, że klasa dziedziczy po specjalizacji szablonu, której argumentem
jest ona sama:

\begin{lstlisting}
template<class Derived>
struct Printable {
    void print() const {
        static_cast<const Derived&>(*this).print_impl();
    }
};

struct Number : Printable<Number> {
    int value = 0;
    void print_impl() const;
};
\end{lstlisting}

Mechanizm \texttt{static\_cast} jest tutaj poprawny pod założeniem,
że obiekt bazowy faktycznie stanowi podobiekt \texttt{Derived}.
CRTP nie zapewnia dynamicznego sprawdzania typu ani wirtualnego
wywołania — poprawność wynika z konstrukcji hierarchii. Wywołanie
\texttt{print\_impl()} może być statycznie związane i często
łatwiej je wbudować w kod wynikowy (\textit{inline}).

Bardziej praktycznym przykładem jest domieszka (\textit{mixin})
dostarczająca wspólną operację, lecz działająca na danych
konkretnej klasy pochodnej:

\begin{lstlisting}
template<class Derived>
struct TranslationMixin {
    constexpr void translate(int dx, int dy) {
        auto& self = static_cast<Derived&>(*this);
        self.x += dx;
        self.y += dy;
    }
};

struct Point : TranslationMixin<Point> {
    int x{}, y{};
    constexpr Point(int a, int b) : x(a), y(b) {}
};

constexpr bool test_translation() {
    Point p{1, 2};
    p.translate(3, 4);
    return p.x == 4 && p.y == 6;
}
static_assert(test_translation());
\end{lstlisting}

Baza \texttt{TranslationMixin<Point>} jest instancjowana z typem
\texttt{Point}, chociaż sam typ \texttt{Point} nie jest jeszcze
kompletny w momencie deklarowania dziedziczenia. Odwołania do składowych
\texttt{x} i \texttt{y} zależą jednak od \texttt{Derived} i są
sprawdzane dopiero podczas instancjowania odpowiedniej funkcji.
W rzeczywistym API warto opisać wymagane składowe koncepcją lub
przenieść operację do jawnie wskazanego interfejsu.

\paragraph{CRTP a wirtualne funkcje.}
Klasyczny polimorfizm dynamiczny pozwala wybierać implementację
na podstawie rzeczywistego typu obiektu przez interfejs klasy bazowej,
także gdy typ jest nieznany podczas kompilacji kodu wywołującego.
CRTP wymaga znajomości konkretnego typu podczas instancjowania.
Nie jest ogólnym zamiennikiem klas abstrakcyjnych. W C++20 prostsze
zastosowania domieszek i interfejsów statycznych można często oprzeć
na \texttt{concepts} i funkcjach wolnostojących.

\section{Policy-based design: wybór zachowania przez typ}
\label{sec:cpp9-policy}

Szablon może przyjąć \emph{politykę}, której metody definiują wariant
algorytmu. Umożliwia to wybór zachowania bez wirtualnych wywołań,
a kompilator może całkowicie usunąć niepotrzebne operacje:

\begin{lstlisting}
#include <limits>
#include <stdexcept>

struct Checked {
    static int divide(int a, int b) {
        if (b == 0) throw std::invalid_argument("zero");
        if (a == std::numeric_limits<int>::min() && b == -1)
            throw std::overflow_error("result outside int");
        return a / b;
    }
};

struct Fast {
    static int divide(int a, int b) { return a / b; }
};

template<class Policy>
struct Calculator {
    int divide(int a, int b) const {
        return Policy::divide(a, b);
    }
};

// Calculator<Checked> checks both invalid integer cases;
// Calculator<Fast> requires a nonzero divisor and no overflow.
\end{lstlisting}

Polityka \texttt{Fast} nie zabezpiecza dzielenia przez zero ani
przepełnienia \texttt{INT\_MIN / -1}; oba przypadki mają
niezdefiniowane zachowanie dla \texttt{int}. \texttt{Checked}
odrzuca je przed dzieleniem. Dobór polityki jest częścią kontraktu API.

W realnym kodzie szablon polityki można ograniczyć koncepcją, a jeśli
liczba wariantów jest niewielka, zwykłe \texttt{if constexpr} może być
łatwiejsze w utrzymaniu niż rozbudowana hierarchia polityk.

\section{Kompletny przykład: dobór typu akumulatora}
\label{sec:cpp9-worked}

Załóżmy, że chcemy obliczać sumę liczb, lecz wybierać typ akumulatora
w zależności od argumentu: dla typów całkowitych \texttt{long long},
a dla zmiennoprzecinkowych \texttt{long double}. Rozwiązanie łączy
\texttt{concepts}, \texttt{type\_traits} i \texttt{conditional\_t}:

\begin{lstlisting}
#include <concepts>
#include <type_traits>
#include <vector>

template<class T>
concept Numeric = std::integral<T> || std::floating_point<T>;

template<Numeric T>
using sum_type_t = std::conditional_t<
    std::integral<T>, long long, long double>;

template<Numeric T>
sum_type_t<T> accumulate_values(const std::vector<T>& xs) {
    sum_type_t<T> result{};
    for (const T& x : xs) {
        result += static_cast<sum_type_t<T>>(x);
    }
    return result;
}

static_assert(std::is_same_v<sum_type_t<int>, long long>);
static_assert(std::is_same_v<sum_type_t<double>, long double>);
\end{lstlisting}

\paragraph{Ograniczenia.}
Przykład wyjaśnia selekcję typu, lecz \emph{nie daje gwarancji
braku przepełnienia}. Nawet \texttt{long long} ma ograniczony zakres,
a suma zmiennoprzecinkowa może akumulować błędy zaokrągleń.
Najpierw należy ustalić wymagania numeryczne, a dopiero potem
projektować odpowiednią metafunkcję.

\section{Czas kompilacji, diagnostyka i praktyczne ograniczenia}

Metaprogramowanie przenosi koszt części pracy z czasu wykonania na
czas kompilacji, ale samo w sobie nie oznacza wydajniejszego programu.
Każda odmienna specjalizacja szablonu może wymagać osobnego
instancjowania i wygenerowania kodu. Powstają wtedy m.in.:
\begin{itemize}
\item dłuższe kompilacje i większe zużycie pamięci kompilatora;
\item \emph{code bloat}, czyli powiększenie plików wynikowych przez wiele
podobnych instancji;
\item skomplikowane komunikaty błędów przy zagnieżdżonych metafunkcjach;
\item problemy z zależnościami między plikami nagłówkowymi;
\item ograniczenia głębokości instancjowania i kosztu ewaluacji stałej.
\end{itemize}

\paragraph{Pliki nagłówkowe i jawne instancjowanie.}
Definicje szablonów zazwyczaj umieszczamy w nagłówkach, ponieważ
kompilator musi je znać przy instancjowaniu z konkretnymi argumentami.
Możliwe jest również \emph{jawne instancjowanie} dla ustalonych typów,
co pozwala kontrolować zestaw generowanych specjalizacji:

\begin{lstlisting}
// In a header:
template<class T> T triple(T x);
extern template int triple<int>(int);

// In one .cpp file (definition must be visible here):
template<class T> T triple(T x) { return x + x + x; }
template int triple<int>(int);
\end{lstlisting}

Ten przykład przedstawia ideę rozdzielenia deklaracji i jawnego
instancjowania; obie oznaczone części należy umieścić w odpowiednich
plikach, a dla wywołań innych typów zapewnić dostęp do definicji.
Warto pamiętać o regule ODR z rozdziału~4.

\section{Jaką technikę wybrać?}

\begin{center}
\small
\begin{tabular}{|p{5.7cm}|p{6.8cm}|}
\hline
\textbf{Potrzeba} & \textbf{Najczęściej odpowiedni mechanizm} \\
\hline
Sprawdzenie wymaganych operacji w nowym API & \texttt{concepts}, \texttt{requires} (C++20) \\
\hline
Obliczenie wartości dla stałych danych & Funkcja \texttt{constexpr} \\
\hline
Wymuszenie ewaluacji w czasie kompilacji & \texttt{consteval} (C++20) \\
\hline
Warunkowy kod zależny od typu & \texttt{if constexpr} \\
\hline
Ustalenie własności typu & \texttt{std::is\_...\_v}, inne \texttt{type\_traits} \\
\hline
Zmiana lub wybór typu wynikowego & \texttt{...\_t}, specjalizacja pomocniczej klasy \\
\hline
Wsparcie biblioteki sprzed C++20 & SFINAE, \texttt{enable\_if}, \texttt{void\_t} \\
\hline
Wieloargumentowe cechy typów & Pakiety typów, \texttt{fold expressions} \\
\hline
Dostęp do pozycji elementów krotki & \texttt{index\_sequence}, \texttt{std::apply} \\
\hline
Wielokrotne użycie zachowania klasowego & CRTP / mixin, rozważ też koncepcje \\
\hline
Wybór wariantu strategii w typie & Szablon polityki (policy-based design) \\
\hline
\end{tabular}
\end{center}

\section{Typowe błędy i zasady projektowe}

\begin{enumerate}
\item \textbf{Mylenie typu z wartością.} \texttt{typename T::value\_type}
      to typ, a \texttt{Trait<T>::value} to wartość; nie są wymienne.
\item \textbf{Brak warunku stopu.} Rekurencyjny szablon klasy może
      prowadzić do nieograniczonego instancjowania i błędów głębokości.
\item \textbf{Założenie, że SFINAE przechwytuje wszystkie błędy.}
      Działa tylko w określonych kontekstach podstawiania.
\item \textbf{Nadużywanie \texttt{conditional\_t}.} Oba argumenty
      typowe trzeba poprawnie sformować, zanim zostanie dokonany wybór.
\item \textbf{Próba wywołania \texttt{declval}.} Można go używać wyłącznie
      w odpowiednich kontekstach niewartościowanych.
\item \textbf{Mylenie \texttt{decltype(x)} z \texttt{decltype((x))}.}
      Dodatkowe nawiasy mogą zmienić wynik na typ referencyjny.
\item \textbf{Mylenie \texttt{constexpr} z obowiązkiem obliczenia
      podczas kompilacji.} Nie każde wywołanie funkcji \texttt{constexpr}
      jest ewaluowane w czasie kompilacji.
\item \textbf{Przenoszenie semantyki do traits.} Typ spełniający
      syntaktyczne wymaganie nie musi spełniać własności biznesowej,
      np. braku utraty precyzji czy bezpiecznej serializacji.
\item \textbf{Niepotrzebny CRTP.} Gdy wystarczy wolna funkcja
      z koncepcją, hierarchia klas często niczego nie upraszcza.
\item \textbf{Przedwczesne metaprogramowanie.} Najprostsze czytelne
      rozwiązanie jest zazwyczaj lepsze od dodatkowej warstwy abstrakcji.
\end{enumerate}

\paragraph{Pytania kontrolne.}
\begin{enumerate}
\item Czym różni się instancjowanie szablonu od wykonania funkcji?
\item Dlaczego \texttt{Factorial<0>} wymaga osobnej specjalizacji?
\item Kiedy \texttt{std::decay\_t<T>} i \texttt{std::remove\_cvref\_t<T>}
      dają różne wyniki?
\item Dlaczego \texttt{std::conditional\_t} nie rozwiązuje wszystkich
      problemów warunkowego tworzenia typów?
\item Co dokładnie sprawdza \texttt{void\_t<decltype(++declval<T\&>())>}?
\item Kiedy niepowodzenie podstawienia jest SFINAE, a kiedy zwykłym
      błędem kompilacji?
\item Dlaczego \texttt{if constexpr} nie zastępuje preprocesora?
\item Czym różnią się \texttt{constexpr} i \texttt{consteval}?
\item W jakich sytuacjach CRTP nie zastąpi funkcji wirtualnych?
\item Kiedy lepiej zastosować \texttt{concepts} niż \texttt{enable\_if}?
\end{enumerate}

\paragraph{Dalsza lektura.}
Warto sięgnąć do dokumentacji \emph{cppreference} dotyczącej
\emph{Templates}, \emph{Type traits} i \emph{SFINAE}.
Pomocne są również hasła \texttt{std::void\_t}, \texttt{std::enable\_if},
\texttt{if constexpr}, \texttt{consteval} oraz
\texttt{std::integer\_sequence}. Do diagnozowania złożonych
błędów szablonów przydają się opcje \texttt{-Wall -Wextra -Wpedantic},
\texttt{-ftemplate-backtrace-limit=0} (GCC) oraz ograniczanie zasięgu
instancjowanego kodu do minimalnego przykładu.


\chapter[Metaprogramowanie w C++]{Zaawansowane metaprogramowanie: szablony, cechy typów i obliczenia podczas kompilacji}
\label{chap:meta-cpp}

\section{Czym jest metaprogramowanie?}

\emph{Metaprogramowanie} oznacza pisanie kodu, który na etapie kompilacji
wyznacza typy, własności typów, wartości lub strukturę później
kompilowanego programu. Szablon nie jest gotową funkcją czy klasą:
stanowi przepis, z którego kompilator tworzy konkretne wersje
dla podanych typów. Nazywamy to \emph{instancjowaniem szablonu}.
Na przykład wywołanie
\texttt{twice(21)} w programie poniżej każe kompilatorowi przygotować
wersję funkcji dla \texttt{int}; program podczas działania wywołuje
już tę wersję. W C++20 do dyspozycji mamy kilka narzędzi:
\begin{itemize}
\item \textbf{szablony klas i funkcji} — generowanie kodu dla różnych typów;
\item \textbf{\texttt{type\_traits}} — obliczanie własności i przekształcanie typów;
\item \textbf{specjalizacje i SFINAE} — wybór poprawnych wersji szablonów;
\item \textbf{\texttt{if constexpr}, \texttt{constexpr}, \texttt{consteval}} —
      wykonywanie obliczeń i selekcja kodu w czasie kompilacji;
\item \textbf{\texttt{concepts}} — czytelne formułowanie wymagań wobec typów;
\item \textbf{CRTP} — wykorzystanie dziedziczenia do polimorfizmu statycznego.
\end{itemize}

Metaprogramowanie ma praktyczne cele: eliminowanie powtarzalnego kodu,
uzyskiwanie lepszych błędów kompilacji, wybór wydajnej implementacji oraz
zapewnienie poprawności określonych operacji \emph{przed uruchomieniem} programu.
Nie oznacza automatycznie zwiększenia wydajności — kosztem mogą być dłuższa
kompilacja, powiększenie kodu maszynowego i trudniejsze komunikaty diagnostyczne.

\begin{lstlisting}
#include <type_traits>

template<class T>
T twice(T x) {
    static_assert(std::is_arithmetic_v<T>);
    return x + x;
}

int a = twice(21);       // T = int
// twice("text");       // static_assert fails
\end{lstlisting}

\paragraph{Związki z innymi rozdziałami.}
Wykorzystamy szablony wariadyczne (rozdział~11),
przenoszenie i referencje (rozdział~6) oraz lambdy (rozdział~9).
Koncepcje szerzej poznamy w rozdziale~13; tutaj pokazujemy
przede wszystkim mechanizmy, na których się opierają.
Wszystkie oznaczone jako poprawne przykłady odnoszą się do \textbf{C++20}.

\section{Instancjowanie szablonów, parametry i zależne nazwy}
\label{sec:cpp9-instantiation}

\subsection{Parametry typowe, nietypowe i szablonowe}

Szablon może przyjmować argumenty będące typami, wartościami stałymi
lub innymi szablonami. Parametr wartości (\textit{non-type template parameter},
NTTP) w C++20 może również mieć odpowiedni typ klasowy spełniający
wymagania tzw. typu strukturalnego.

\begin{lstlisting}
#include <array>
#include <cstddef>

template<class T, std::size_t N>
struct Buffer {
    std::array<T, N> data{};
};

Buffer<double, 4> samples; // T=double, N=4
\end{lstlisting}

\texttt{Buffer<double,4>} i \texttt{Buffer<double,8>} są różnymi typami.
Parametr \texttt{N} jest znany podczas kompilacji i może wystąpić tam,
gdzie język wymaga wyrażenia stałego, np. jako rozmiar \texttt{std::array}.

\subsection{Parametr będący obiektem klasy (C++20)}

Od C++20 odpowiednie \emph{typy strukturalne} mogą występować jako
parametry szablonu będące wartościami. Umożliwia to np. identyfikowanie
szablonu przez napis znany podczas kompilacji:

\begin{lstlisting}
#include <cstddef>

template<std::size_t N>
struct fixed_string {
    char data[N]{};
    constexpr fixed_string(const char (&s)[N]) {
        for (std::size_t i = 0; i < N; ++i)
            data[i] = s[i];
    }
};

template<fixed_string Name>
struct NamedCommand {
    static constexpr auto name = Name;
};

static_assert(NamedCommand<"run">::name.data[0] == 'r');
\end{lstlisting}

To przykład zaawansowanej techniki; zwykły obiekt \texttt{std::string}
nie jest w C++20 strukturalnym argumentem szablonu. Rozmiar i
zawartość napisu wpływają tutaj na tożsamość specjalizacji szablonu.

\subsection{Nazwy zależne: \texttt{typename} i \texttt{template}}

W definicji szablonu kompilator nie zawsze wie, czy nazwa zależna od
parametru jest typem, czy wartością. W wielu kontekstach musimy to
jednoznacznie zaznaczyć:

\begin{lstlisting}
#include <vector>

template<class Container>
auto first_copy(const Container& c) {
    typename Container::value_type x = c.front();
    return x;
}

template<class T>
struct Wrapper {
    template<class U>
    static constexpr bool same = false;
};

template<class T>
constexpr bool use_member_template() {
    return Wrapper<T>::template same<int>;
}
\end{lstlisting}

Słowo \texttt{typename} wskazuje, że zależna nazwa członkowska
jest typem. Zapis \texttt{::template same<int>} informuje,
że po operatorze \texttt{::} występuje szablon składowy, a znaki
\texttt{<} i \texttt{>} są delimitatorami argumentów szablonu.
Dokładne reguły zależą od kontekstu; obu słów nie należy dopisywać
mechanicznie wszędzie.

\section{Metafunkcje i \texttt{std::integral\_constant}}
\label{sec:cpp9-integral-constant}

Poniższa silnia obliczana przez szablon pokazuje starszy styl
metaprogramowania. Gdy wynikiem ma być zwykła liczba, funkcja
\texttt{constexpr} jest dziś zazwyczaj prostsza. Szablony klas
pozostają potrzebne zwłaszcza wtedy, gdy wynikiem obliczenia ma być
\emph{typ}.

Klasyczne metaprogramowanie szablonowe opierało się na klasach, które
\emph{przenosiły wyniki obliczenia w swoich składowych statycznych}.
Dla przykładu silnię można wyznaczyć przez rekurencyjne instancjowanie
szablonów:

\begin{lstlisting}
template<unsigned N>
struct Factorial {
    static constexpr unsigned value =
        N * Factorial<N - 1>::value;
};

template<>
struct Factorial<0> {
    static constexpr unsigned value = 1;
};

static_assert(Factorial<5>::value == 120);
\end{lstlisting}

Specjalizacja dla zera pełni rolę warunku stopu. Bez niej kompilator
próbowałby tworzyć kolejne instancje, aż osiągnąłby limit głębokości.
Obliczenie ma charakter rekursji \emph{w strukturze typów}, a nie
rekursji wywołań funkcji wykonywanych w czasie działania programu.

Standardowym fundamentem takich metafunkcji jest
\texttt{std::integral\_constant<T,v>}:

\begin{lstlisting}
#include <type_traits>

using TrueType = std::integral_constant<bool, true>;
static_assert(TrueType::value);
static_assert(std::bool_constant<true>::value);
static_assert(std::true_type::value);
static_assert(!std::false_type::value);

template<unsigned N>
struct PowerOfTwo :
    std::integral_constant<unsigned, 2 * PowerOfTwo<N-1>::value> {};

template<>
struct PowerOfTwo<0> : std::integral_constant<unsigned, 1> {};

static_assert(PowerOfTwo<5>::value == 32);
\end{lstlisting}

Predykaty takie jak \texttt{std::is\_pointer} udostępniają wartość
\texttt{::value}. Od C++17 można używać krótszej wersji
z przyrostkiem \texttt{\_v}, na przykład:
\begin{center}
\texttt{std::is\_pointer\_v<T>}.
\end{center}

\subsection{Metafunkcja szablonowa czy funkcja \texttt{constexpr}?}

\begin{lstlisting}
constexpr unsigned factorial(unsigned n) {
    unsigned result = 1;
    for (unsigned i = 2; i <= n; ++i)
        result *= i;
    return result;
}

static_assert(factorial(5) == 120);
unsigned x = 6;
unsigned y = factorial(x); // allowed at run time
\end{lstlisting}

Jeśli wynik jest zwykłą wartością, funkcja \texttt{constexpr} jest
najczęściej prostsza niż rekurencyjna metafunkcja. Klasyczne metafunkcje
pozostają użyteczne, gdy \emph{wynikiem jest typ}, gdy korzystamy ze
specjalizacji lub interoperujemy ze starszym API.

\section{Specjalizacje: pełne i częściowe}
\label{sec:cpp9-specialization}

Szablon podstawowy definiuje przypadek ogólny. \emph{Specjalizacja pełna}
dotyczy konkretnego zestawu argumentów, a \emph{specjalizacja częściowa}
obsługuje pewną ich rodzinę:

\begin{lstlisting}
#include <type_traits>

template<class T>
struct is_raw_pointer : std::false_type {};

template<class T>
struct is_raw_pointer<T*> : std::true_type {};

static_assert(!is_raw_pointer<int>::value);
static_assert(is_raw_pointer<double*>::value);
static_assert(!is_raw_pointer<int* const>::value);
\end{lstlisting}

Ostatni wynik nie jest pomyłką: \texttt{int* const} jest
\emph{stałym wskaźnikiem}, nie typem postaci \texttt{T*}.
Jeśli chcemy ignorować kwalifikatory zewnętrzne, możemy najpierw
użyć \texttt{std::remove\_cv\_t<T>}.

\begin{lstlisting}
template<class T>
inline constexpr bool is_any_raw_pointer_v =
    is_raw_pointer<std::remove_cv_t<T>>::value;

static_assert(is_any_raw_pointer_v<int* const>);
\end{lstlisting}

\paragraph{Ważne ograniczenie.}
Dla \textbf{szablonów funkcji nie ma specjalizacji częściowej}.
W takich przypadkach wykorzystujemy przeciążanie, koncepcje,
\texttt{if constexpr} lub specjalizację pomocniczego szablonu klasy.
Pełna specjalizacja funkcji istnieje, ale jej interakcja z przeciążeniami
bywa zaskakująca; w nowym kodzie zwykle preferujemy przeciążenia.

\section{\texttt{type\_traits}: pytania o typy}
\label{sec:cpp9-traits}

Biblioteka \texttt{<type\_traits>} jest zbiorem standardowych metafunkcji.
Dzielimy je na \emph{predykaty} (wynik logiczny) i \emph{transformacje}
(wynik będący typem). Ich składnia opiera się na ustalonych konwencjach:
\begin{itemize}
\item \texttt{Trait<T>::value} lub \texttt{Trait\_v<T>} dla predykatów;
\item \texttt{typename Trait<T>::type} lub \texttt{Trait\_t<T>}
      dla transformacji.
\end{itemize}

\begin{lstlisting}
#include <type_traits>
#include <string>

static_assert(std::is_integral_v<int>);
static_assert(std::is_floating_point_v<double>);
static_assert(std::is_same_v<int, int>);
static_assert(std::is_pointer_v<const char*>);
static_assert(std::is_constructible_v<std::string, const char*>);
static_assert(std::is_move_constructible_v<std::string>);

using A = std::remove_const_t<const int>; // int
using B = std::remove_reference_t<int&>; // int
using C = std::remove_cvref_t<const int&>; // int (C++20)
using D = std::add_pointer_t<int>;        // int*
static_assert(std::is_same_v<A, B>);
static_assert(std::is_same_v<C, int>);
static_assert(std::is_same_v<D, int*>);
\end{lstlisting}

\begin{center}
\small
\begin{tabular}{|>{\raggedright\arraybackslash}p{5.5cm}|>{\raggedright\arraybackslash}p{7.0cm}|}
\hline
\textbf{Cecha} & \textbf{Co sprawdza lub wyznacza?} \\
\hline
\texttt{is\_same\_v<A,B>} & Identyczność typów (wraz z istotnymi kwalifikatorami). \\
\hline
\texttt{is\_convertible\_v<A,B>} & Dostępność konwersji niejawnej \texttt{A} do \texttt{B}. \\
\hline
\texttt{is\_constructible\_v<T,Args...>} & Możliwość skonstruowania \texttt{T} z argumentów zadanych typów. \\
\hline
\texttt{is\_nothrow\_move\_}\newline\texttt{constructible\_v<T>} & Deklarowana bezwyjątkowość konstrukcji przenoszącej. \\
\hline
\texttt{is\_trivially\_copyable\_v<T>} & Własność pozwalająca na określone operacje z reprezentacją bajtową. \\
\hline
\texttt{remove\_cvref\_t<T>} & Usunięcie referencji i kwalifikatorów \texttt{const}/\texttt{volatile} z poziomu wierzchniego. \\
\hline
\texttt{common\_type\_t<A,B>} & Wspólny typ wynikowy dla danych typów (o ile istnieje). \\
\hline
\texttt{invoke\_result\_t<F,Args...>} & Typ wyniku wywołania poprzez \texttt{std::invoke}. \\
\hline
\end{tabular}
\end{center}

\paragraph{\texttt{remove\_cvref} a \texttt{decay}.}
\texttt{std::decay\_t<T>} oprócz usunięcia referencji i zewnętrznych
kwalifikatorów może przekształcić tablicę w wskaźnik oraz typ funkcji
w wskaźnik do funkcji. Nie są to zatem synonimy.

\begin{lstlisting}
using X = std::remove_cvref_t<int(&)[3]>; // int[3]
using Y = std::decay_t<int(&)[3]>;        // int*
static_assert(std::is_same_v<Y, int*>);
\end{lstlisting}

\paragraph{Czego predykat nie gwarantuje?}
\texttt{std::is\_move\_constructible\_v<T>} może być prawdziwe również
wtedy, gdy wywołanie z argumentem rvalue używa konstruktora kopiującego.
Podobnie \texttt{is\_trivially\_copyable} nie oznacza, że każdy odczyt
bajtów, rzutowanie lub serializacja do pliku będzie poprawna bądź
przenośna. Cechy typów opisują wybrane własności języka, nie
pełny kontrakt semantyczny.

\subsection{Składanie predykatów z zachowaniem krótkiego obliczania}

\texttt{std::conjunction}, \texttt{std::disjunction} i
\texttt{std::negation} łączą predykaty, które mają składową
\texttt{::value}. Dwa pierwsze narzędzia zatrzymują odczytywanie \texttt{::value}
z kolejnych predykatów, kiedy wynik jest już znany. Trzeba jednak
poprawnie sformować same argumenty typowe tych metafunkcji.

\begin{lstlisting}
#include <type_traits>

template<class T>
using SmallSignedInteger = std::conjunction<
    std::is_integral<T>,
    std::is_signed<T>,
    std::bool_constant<(sizeof(T) <= sizeof(int))>>;

static_assert(SmallSignedInteger<short>::value);
static_assert(!SmallSignedInteger<unsigned>::value);
\end{lstlisting}

Dla prostych stałych logicznych zapis
\texttt{(cond1 \&\& cond2)} bywa czytelniejszy.
\texttt{conjunction} jest użyteczne przede wszystkim wtedy, gdy
poszczególne warunki są odrębnymi metafunkcjami, a ich instancjowanie
może być kosztowne lub niepoprawne bez wcześniejszego warunku.

\section{Metafunkcje zwracające typy i leniwy wybór}
\label{sec:cpp9-transforms}

Tworząc własną transformację typu, zwyczajowo zapisujemy wynik jako
\texttt{using type = ...} i udostępniamy wygodny alias:

\begin{lstlisting}
template<class T>
struct storage_type {
    using type = std::remove_cvref_t<T>;
};

template<class T>
using storage_type_t = typename storage_type<T>::type;

static_assert(std::is_same_v<storage_type_t<const long&>, long>);
\end{lstlisting}

Ważnym narzędziem jest \texttt{std::conditional\_t<B,T,F>}, wybierające
jeden z dwóch typów w zależności od znanej w kompilacji wartości logicznej:

\begin{lstlisting}
template<class T>
using accumulator_t = std::conditional_t<
    std::is_integral_v<T>, long long, long double>;

static_assert(std::is_same_v<accumulator_t<int>, long long>);
static_assert(std::is_same_v<accumulator_t<float>, long double>);
\end{lstlisting}

\paragraph{Pułapka eager instantiation.}
\texttt{conditional\_t} wybiera typ spośród \emph{już utworzonych
argumentów typowych}. Nie chroni przed błędami powstającymi podczas
formowania któregoś z nich:

\begin{lstlisting}
// Incorrect for int: typename T::value_type is invalid.
// template<class T>
// using bad = std::conditional_t<
//     std::is_integral_v<T>, T, typename T::value_type>;
\end{lstlisting}

Aby uzyskać prawdziwie warunkowe tworzenie typu, można zastosować
specjalizację pomocniczej klasy:

\begin{lstlisting}
template<class T, bool = std::is_integral_v<T>>
struct value_or_self {
    using type = typename T::value_type;
};

template<class T>
struct value_or_self<T, true> {
    using type = T;
};

template<class T>
using value_or_self_t = typename value_or_self<T>::type;

static_assert(std::is_same_v<value_or_self_t<int>, int>);
\end{lstlisting}

Podstawowa definicja zawiera nazwę zależną, która jest sprawdzana przy
instancjowaniu odpowiedniego przypadku; nie jest potrzebna dla \texttt{int}.

\section{\texttt{decltype}, \texttt{declval} i konteksty bez obliczania}

\begin{center}\small
\begin{tabularx}{\textwidth}{@{}p{0.32\textwidth}X@{}}
\toprule
Narzędzie & Za co odpowiada \\
\midrule
\code{decltype(expr)} & Ustala typ wyrażenia bez wykonywania go; szczególna reguła dotyczy nieujętej w nawiasy nazwy. \\
\code{decltype((x))} & Dla lvalue \code{x} zwraca typ referencyjny, np. \texttt{T\&}. \\
\code{std::declval<T>()} & Udostępnia fikcyjne wyrażenie typu \code{T} wyłącznie w kontekście nieobliczanym. \\
\code{std::void_t<...>} & Zamienia poprawne podstawienia typów na \code{void}; przydatne w starszych technikach wykrywania. \\
\code{if constexpr} & W szablonie wybiera gałąź podczas kompilacji i odrzuca niepotrzebną gałąź po instancjacji. \\
\bottomrule
\end{tabularx}
\end{center}
\label{sec:cpp9-decltype}

\texttt{decltype(expr)} pozwala odczytać typ wyrażenia bez wykonywania
samego wyrażenia. Subtelność stanowi różnica między nieujętym w nawiasy
identyfikatorem a ogólnym wyrażeniem:

\begin{lstlisting}
int x = 5;
static_assert(std::is_same_v<decltype(x), int>);
static_assert(std::is_same_v<decltype((x)), int&>);
\end{lstlisting}

Wyrażenie \texttt{(x)} jest \textit{lvalue}, więc \texttt{decltype((x))}
zwraca \texttt{int\&}. Reguła ta ma bezpośredni związek z kategoriami
wartości (rozdział~6).

\texttt{std::declval<T>()} pozwala użyć argumentu typu \texttt{T}
w \emph{kontekście niewartościowanym}, nawet jeśli nie możemy utworzyć
jego obiektu. Nie wolno rzeczywiście wykonywać wywołania
\texttt{declval}.

\begin{lstlisting}
#include <utility>
#include <type_traits>

template<class T>
using dereference_t = decltype(*std::declval<T&>());

static_assert(std::is_same_v<dereference_t<int*>, int&>);

struct NonDefault {
    NonDefault() = delete;
    int method() const;
};

using Result = decltype(std::declval<const NonDefault&>().method());
static_assert(std::is_same_v<Result, int>);
\end{lstlisting}

Typ wyniku wywołania można również wyznaczyć za pomocą gotowej metafunkcji:
\begin{center}
\texttt{std::invoke\_result\_t<F, Args...>}.
\end{center}

\section{SFINAE: błędne podstawienie to nie zawsze błąd kompilacji}
\label{sec:cpp9-sfinae}

Ta część pomaga czytać biblioteki napisane przed C++20.
W nowym kodzie wymagania wobec typów zwykle prościej wyrazić
przez \texttt{requires}, opisane w rozdziale~13.

\textbf{SFINAE} (\textit{Substitution Failure Is Not An Error}) mówi,
że niektóre błędy powstałe podczas podstawiania argumentów do szablonu
\emph{nie kończą kompilacji}: dana specjalizacja zostaje wyłączona
z zestawu kandydatów lub odpowiednia specjalizacja częściowa nie pasuje.
Dotyczy to określonych kontekstów bezpośredniego podstawiania;
\textbf{nie} jest mechanizmem przechwytującym dowolne błędy w ciele funkcji
ani skutki uboczne instancjowania innych typów.

\subsection{\texttt{std::enable\_if\_t}}

\begin{lstlisting}
#include <type_traits>

template<class T>
constexpr std::enable_if_t<std::is_integral_v<T>, T>
doubled(T x) {
    return x + x;
}

static_assert(doubled(3) == 6);
// doubled(2.5); // no matching function
\end{lstlisting}

Dla typu \texttt{double} nie istnieje \texttt{type} w
\texttt{std::enable\_if<false,double>}; funkcja zostaje wyłączona
z odpowiedniego zbioru kandydatów. Do konstruktorów, operatorów i
różnych wariantów przeciążeń istnieją odmienne miejsca umieszczania
\texttt{enable\_if}; łatwo przy tym uzyskać przypadkowe powtórne
deklaracje lub nieczytelne błędy.

\subsection{\texttt{void\_t} — detection idiom}

Od C++17 wygodnie wykrywać istnienie nazw i poprawność wyrażeń przy
użyciu \texttt{std::void\_t}. Jeżeli wszystkie podane argumenty
są prawidłowymi typami, rezultatem jest po prostu \texttt{void}:

\begin{lstlisting}
#include <type_traits>
#include <vector>

template<class, class = void>
struct has_value_type : std::false_type {};

template<class T>
struct has_value_type<T, std::void_t<typename T::value_type>>
    : std::true_type {};

static_assert(has_value_type<std::vector<int>>::value);
static_assert(!has_value_type<int>::value);
\end{lstlisting}

W podstawieniu \texttt{T=int} wyrażenie \texttt{int::value\_type}
nie ma sensu; specjalizacja częściowa nie jest więc wybierana.
Analogicznie można badać poprawność wyrażenia z \texttt{decltype}:

\begin{lstlisting}
#include <type_traits>
#include <utility>

template<class, class = void>
struct has_preincrement : std::false_type {};

template<class T>
struct has_preincrement<T,
    std::void_t<decltype(++std::declval<T&>())>>
    : std::true_type {};

static_assert(has_preincrement<int>::value);
\end{lstlisting}

\subsection{SFINAE kontra \texttt{requires} (C++20)}

Konstrukcję wykrywającą operację możemy wyrazić bezpośrednio za pomocą
koncepcji z rozdziału~13:

\begin{lstlisting}
#include <concepts>

template<class T>
concept Preincrementable = requires(T& x) {
    { ++x } -> std::same_as<T&>;
};

template<Preincrementable T>
void advance_one(T& x) { ++x; }
\end{lstlisting}

Metafunkcja \texttt{has\_preincrement<T>} sprawdza jedynie poprawność
operacji \texttt{++x}. Z kolei koncepcja \texttt{Preincrementable}
dodatkowo wymaga, aby jej \emph{typ wyniku} odpowiadał dokładnie \texttt{T\&}.
Różnica ta jest zamierzona: ograniczenia należy definiować zgodnie z
kontraktem, a nie tylko według dostępności pewnego wyrażenia.

\section{\texttt{if constexpr}: wybór implementacji na podstawie typu}
\label{sec:cpp9-ifconstexpr}

Od C++17 często zamiast rekurencyjnych specjalizacji wybieramy gałąź
wewnątrz pojedynczego szablonu. Niebrana gałąź zależna od argumentu
szablonu nie jest instancjowana dla danego zestawu argumentów:

\begin{lstlisting}
#include <type_traits>

template<class T>
constexpr auto number_or_value(const T& x) {
    if constexpr (std::is_arithmetic_v<T>)
        return x;
    else
        return x.value();
}

struct Box {
    constexpr int value() const { return 7; }
};
static_assert(number_or_value(3) == 3);
static_assert(number_or_value(Box{}) == 7);
\end{lstlisting}

\paragraph{Dlaczego nie zwykłe \texttt{if}?}
W zwykłym \texttt{if} kompilator musi przeanalizować obie gałęzie
pod kątem poprawności dla instancjowanego kodu. Dla \texttt{int}
wyrażenie \texttt{x.value()} nie jest poprawne, a sam warunek
\texttt{std::is\_arithmetic\_v<int>} tego nie zmieni.

\paragraph{Ważna subtelność.}
\texttt{if constexpr(false)} nie jest zamiennikiem
\texttt{\#if 0}: poza kontekstem szablonowym odrzucona gałąź nadal
podlega standardowym kontrolom poprawności. Ponadto mechanizm nie
usprawiedliwia konstrukcji zawsze niepoprawnych niezależnie od
argumentów szablonu.

\section{Pakiety typów, sekwencje indeksów i fold expressions}
\label{sec:cpp9-type-lists}

Metaprogramowanie dobrze współpracuje z pakietami parametrów
(rozdział~11). Zamiast tworzyć klasyczne rekurencyjne listy typów,
często możemy używać rozwinięcia pakietu i wyrażenia zwijającego.

\begin{lstlisting}
#include <type_traits>

template<class... Ts>
struct type_list {};

template<class T, class... Ts>
inline constexpr bool contains_v =
    (std::is_same_v<T, Ts> || ...); // empty pack -> false

static_assert(contains_v<int, char, int, double>);
static_assert(!contains_v<long, char, int, double>);
static_assert(!contains_v<int>);
\end{lstlisting}

Zamiast tworzyć specjalizacje dla każdego elementu pakietu, korzystamy
z jednego \texttt{fold expression}. Gdy potrzebna jest reprezentacja
indeksów \texttt{0,1,...,N-1}, biblioteka udostępnia
\texttt{std::make\_index\_sequence<N>}:

\begin{lstlisting}
#include <tuple>
#include <utility>
#include <cstddef>

template<class Tuple, std::size_t... I>
constexpr auto reversed_impl(const Tuple& t,
                             std::index_sequence<I...>) {
    return std::make_tuple(std::get<sizeof...(I)-1-I>(t)...);
}

template<class... Ts>
constexpr auto reversed(const std::tuple<Ts...>& t) {
    return reversed_impl(t, std::index_sequence_for<Ts...>{});
}

static_assert(reversed(std::tuple{1, 2, 3}) ==
              std::tuple{3, 2, 1});
\end{lstlisting}

\paragraph{Czy zawsze potrzebujemy indeksów?}
Nie. W przypadku prostego wykonania funkcji na elementach krotki
w C++17 często wystarczy \texttt{std::apply}; sekwencja indeksów
jest potrzebna zwłaszcza wtedy, gdy chcemy \emph{przekształcić pozycje}
lub mieć indywidualny dostęp do każdego indeksu na etapie kompilacji.

\section{\texttt{constexpr}, \texttt{consteval} i obliczenia na danych}
\label{sec:cpp9-constant-eval}

Szablony nie są jedynym sposobem wykonania pracy w trakcie kompilacji.
Od C++11 rozwijano możliwości \texttt{constexpr}, a w C++20 mamy
również \texttt{consteval}. Różnica jest istotna:
\begin{itemize}
\item funkcja \texttt{constexpr} \emph{może} uczestniczyć w obliczeniach
stałych, a w dozwolonych wywołaniach działa też podczas wykonania;
\item wywołanie funkcji \texttt{consteval} musi spełniać reguły
właściwe obliczaniu w czasie kompilacji (\textit{immediate function}).
\end{itemize}

\begin{lstlisting}
constexpr int square(int x) { return x * x; }
consteval int cube_ct(int x) { return x * x * x; }

static_assert(square(4) == 16);
static_assert(cube_ct(3) == 27);

int n = 4;
int a = square(n); // runtime evaluation possible
// int b = cube_ct(n); // error: n is not constant
\end{lstlisting}

Przykładem praktycznym jest haszowanie \texttt{std::string\_view}
w stałym wyrażeniu. Poniższy wariant FNV-1a nadaje się do zbudowania
stałej identyfikującej np. nazwę polecenia. Nie jest to hash
kryptograficzny i nie wyklucza kolizji.

\begin{lstlisting}
#include <cstdint>
#include <string_view>

consteval std::uint64_t fnv1a(std::string_view text) {
    std::uint64_t h = 14695981039346656037ull;
    for (unsigned char c : text) {
        h ^= c;
        h *= 1099511628211ull;
    }
    return h;
}

static_assert(fnv1a("abc") == 0xe71fa2190541574bull);
\end{lstlisting}

\paragraph{\texttt{constinit} i \texttt{std::is\_constant\_evaluated}.}
\texttt{constinit} zapewnia statyczną inicjalizację odpowiedniej
zmiennej, lecz sam nie czyni jej stałą. Funkcja
\texttt{std::is\_constant\_evaluated()} (C++20) pozwala w
pewnych przypadkach rozpoznać, czy bieżąca ewaluacja ma charakter
stały; nadużywanie jej może jednak zaciemniać semantykę funkcji.

\section{CRTP: polimorfizm statyczny}
\label{sec:cpp9-crtp}

\textbf{CRTP} (\textit{Curiously Recurring Template Pattern}) polega na
tym, że klasa dziedziczy po specjalizacji szablonu, której argumentem
jest ona sama:

\begin{lstlisting}
template<class Derived>
struct Printable {
    void print() const {
        static_cast<const Derived&>(*this).print_impl();
    }
};

struct Number : Printable<Number> {
    int value = 0;
    void print_impl() const;
};
\end{lstlisting}

Mechanizm \texttt{static\_cast} jest tutaj poprawny pod założeniem,
że obiekt bazowy faktycznie stanowi podobiekt \texttt{Derived}.
CRTP nie zapewnia dynamicznego sprawdzania typu ani wirtualnego
wywołania — poprawność wynika z konstrukcji hierarchii. Wywołanie
\texttt{print\_impl()} może być statycznie związane i często
łatwiej je wbudować w kod wynikowy (\textit{inline}).

Bardziej praktycznym przykładem jest domieszka (\textit{mixin})
dostarczająca wspólną operację, lecz działająca na danych
konkretnej klasy pochodnej:

\begin{lstlisting}
template<class Derived>
struct TranslationMixin {
    constexpr void translate(int dx, int dy) {
        auto& self = static_cast<Derived&>(*this);
        self.x += dx;
        self.y += dy;
    }
};

struct Point : TranslationMixin<Point> {
    int x{}, y{};
    constexpr Point(int a, int b) : x(a), y(b) {}
};

constexpr bool test_translation() {
    Point p{1, 2};
    p.translate(3, 4);
    return p.x == 4 && p.y == 6;
}
static_assert(test_translation());
\end{lstlisting}

Baza \texttt{TranslationMixin<Point>} jest instancjowana z typem
\texttt{Point}, chociaż sam typ \texttt{Point} nie jest jeszcze
kompletny w momencie deklarowania dziedziczenia. Odwołania do składowych
\texttt{x} i \texttt{y} zależą jednak od \texttt{Derived} i są
sprawdzane dopiero podczas instancjowania odpowiedniej funkcji.
W rzeczywistym API warto opisać wymagane składowe koncepcją lub
przenieść operację do jawnie wskazanego interfejsu.

\paragraph{CRTP a wirtualne funkcje.}
Klasyczny polimorfizm dynamiczny pozwala wybierać implementację
na podstawie rzeczywistego typu obiektu przez interfejs klasy bazowej,
także gdy typ jest nieznany podczas kompilacji kodu wywołującego.
CRTP wymaga znajomości konkretnego typu podczas instancjowania.
Nie jest ogólnym zamiennikiem klas abstrakcyjnych. W C++20 prostsze
zastosowania domieszek i interfejsów statycznych można często oprzeć
na \texttt{concepts} i funkcjach wolnostojących.

\section{Policy-based design: wybór zachowania przez typ}
\label{sec:cpp9-policy}

Szablon może przyjąć \emph{politykę}, której metody definiują wariant
algorytmu. Umożliwia to wybór zachowania bez wirtualnych wywołań,
a kompilator może całkowicie usunąć niepotrzebne operacje:

\begin{lstlisting}
#include <limits>
#include <stdexcept>

struct Checked {
    static int divide(int a, int b) {
        if (b == 0) throw std::invalid_argument("zero");
        if (a == std::numeric_limits<int>::min() && b == -1)
            throw std::overflow_error("result outside int");
        return a / b;
    }
};

struct Fast {
    static int divide(int a, int b) { return a / b; }
};

template<class Policy>
struct Calculator {
    int divide(int a, int b) const {
        return Policy::divide(a, b);
    }
};

// Calculator<Checked> checks both invalid integer cases;
// Calculator<Fast> requires a nonzero divisor and no overflow.
\end{lstlisting}

Polityka \texttt{Fast} nie zabezpiecza dzielenia przez zero ani
przepełnienia \texttt{INT\_MIN / -1}; oba przypadki mają
niezdefiniowane zachowanie dla \texttt{int}. \texttt{Checked}
odrzuca je przed dzieleniem. Dobór polityki jest częścią kontraktu API.

W realnym kodzie szablon polityki można ograniczyć koncepcją, a jeśli
liczba wariantów jest niewielka, zwykłe \texttt{if constexpr} może być
łatwiejsze w utrzymaniu niż rozbudowana hierarchia polityk.

\section{Kompletny przykład: dobór typu akumulatora}
\label{sec:cpp9-worked}

Załóżmy, że chcemy obliczać sumę liczb, lecz wybierać typ akumulatora
w zależności od argumentu: dla typów całkowitych \texttt{long long},
a dla zmiennoprzecinkowych \texttt{long double}. Rozwiązanie łączy
\texttt{concepts}, \texttt{type\_traits} i \texttt{conditional\_t}:

\begin{lstlisting}
#include <concepts>
#include <type_traits>
#include <vector>

template<class T>
concept Numeric = std::integral<T> || std::floating_point<T>;

template<Numeric T>
using sum_type_t = std::conditional_t<
    std::integral<T>, long long, long double>;

template<Numeric T>
sum_type_t<T> accumulate_values(const std::vector<T>& xs) {
    sum_type_t<T> result{};
    for (const T& x : xs) {
        result += static_cast<sum_type_t<T>>(x);
    }
    return result;
}

static_assert(std::is_same_v<sum_type_t<int>, long long>);
static_assert(std::is_same_v<sum_type_t<double>, long double>);
\end{lstlisting}

\paragraph{Ograniczenia.}
Przykład wyjaśnia selekcję typu, lecz \emph{nie daje gwarancji
braku przepełnienia}. Nawet \texttt{long long} ma ograniczony zakres,
a suma zmiennoprzecinkowa może akumulować błędy zaokrągleń.
Najpierw należy ustalić wymagania numeryczne, a dopiero potem
projektować odpowiednią metafunkcję.

\section{Czas kompilacji, diagnostyka i praktyczne ograniczenia}

Metaprogramowanie przenosi koszt części pracy z czasu wykonania na
czas kompilacji, ale samo w sobie nie oznacza wydajniejszego programu.
Każda odmienna specjalizacja szablonu może wymagać osobnego
instancjowania i wygenerowania kodu. Powstają wtedy m.in.:
\begin{itemize}
\item dłuższe kompilacje i większe zużycie pamięci kompilatora;
\item \emph{code bloat}, czyli powiększenie plików wynikowych przez wiele
podobnych instancji;
\item skomplikowane komunikaty błędów przy zagnieżdżonych metafunkcjach;
\item problemy z zależnościami między plikami nagłówkowymi;
\item ograniczenia głębokości instancjowania i kosztu ewaluacji stałej.
\end{itemize}

\paragraph{Pliki nagłówkowe i jawne instancjowanie.}
Definicje szablonów zazwyczaj umieszczamy w nagłówkach, ponieważ
kompilator musi je znać przy instancjowaniu z konkretnymi argumentami.
Możliwe jest również \emph{jawne instancjowanie} dla ustalonych typów,
co pozwala kontrolować zestaw generowanych specjalizacji:

\begin{lstlisting}
// In a header:
template<class T> T triple(T x);
extern template int triple<int>(int);

// In one .cpp file (definition must be visible here):
template<class T> T triple(T x) { return x + x + x; }
template int triple<int>(int);
\end{lstlisting}

Ten przykład przedstawia ideę rozdzielenia deklaracji i jawnego
instancjowania; obie oznaczone części należy umieścić w odpowiednich
plikach, a dla wywołań innych typów zapewnić dostęp do definicji.
Warto pamiętać o regule ODR z rozdziału~4.

\section{Jaką technikę wybrać?}

\begin{center}
\small
\begin{tabular}{|p{5.7cm}|p{6.8cm}|}
\hline
\textbf{Potrzeba} & \textbf{Najczęściej odpowiedni mechanizm} \\
\hline
Sprawdzenie wymaganych operacji w nowym API & \texttt{concepts}, \texttt{requires} (C++20) \\
\hline
Obliczenie wartości dla stałych danych & Funkcja \texttt{constexpr} \\
\hline
Wymuszenie ewaluacji w czasie kompilacji & \texttt{consteval} (C++20) \\
\hline
Warunkowy kod zależny od typu & \texttt{if constexpr} \\
\hline
Ustalenie własności typu & \texttt{std::is\_...\_v}, inne \texttt{type\_traits} \\
\hline
Zmiana lub wybór typu wynikowego & \texttt{...\_t}, specjalizacja pomocniczej klasy \\
\hline
Wsparcie biblioteki sprzed C++20 & SFINAE, \texttt{enable\_if}, \texttt{void\_t} \\
\hline
Wieloargumentowe cechy typów & Pakiety typów, \texttt{fold expressions} \\
\hline
Dostęp do pozycji elementów krotki & \texttt{index\_sequence}, \texttt{std::apply} \\
\hline
Wielokrotne użycie zachowania klasowego & CRTP / mixin, rozważ też koncepcje \\
\hline
Wybór wariantu strategii w typie & Szablon polityki (policy-based design) \\
\hline
\end{tabular}
\end{center}

\section{Typowe błędy i zasady projektowe}

\begin{enumerate}
\item \textbf{Mylenie typu z wartością.} \texttt{typename T::value\_type}
      to typ, a \texttt{Trait<T>::value} to wartość; nie są wymienne.
\item \textbf{Brak warunku stopu.} Rekurencyjny szablon klasy może
      prowadzić do nieograniczonego instancjowania i błędów głębokości.
\item \textbf{Założenie, że SFINAE przechwytuje wszystkie błędy.}
      Działa tylko w określonych kontekstach podstawiania.
\item \textbf{Nadużywanie \texttt{conditional\_t}.} Oba argumenty
      typowe trzeba poprawnie sformować, zanim zostanie dokonany wybór.
\item \textbf{Próba wywołania \texttt{declval}.} Można go używać wyłącznie
      w odpowiednich kontekstach niewartościowanych.
\item \textbf{Mylenie \texttt{decltype(x)} z \texttt{decltype((x))}.}
      Dodatkowe nawiasy mogą zmienić wynik na typ referencyjny.
\item \textbf{Mylenie \texttt{constexpr} z obowiązkiem obliczenia
      podczas kompilacji.} Nie każde wywołanie funkcji \texttt{constexpr}
      jest ewaluowane w czasie kompilacji.
\item \textbf{Przenoszenie semantyki do traits.} Typ spełniający
      syntaktyczne wymaganie nie musi spełniać własności biznesowej,
      np. braku utraty precyzji czy bezpiecznej serializacji.
\item \textbf{Niepotrzebny CRTP.} Gdy wystarczy wolna funkcja
      z koncepcją, hierarchia klas często niczego nie upraszcza.
\item \textbf{Przedwczesne metaprogramowanie.} Najprostsze czytelne
      rozwiązanie jest zazwyczaj lepsze od dodatkowej warstwy abstrakcji.
\end{enumerate}

\paragraph{Pytania kontrolne.}
\begin{enumerate}
\item Czym różni się instancjowanie szablonu od wykonania funkcji?
\item Dlaczego \texttt{Factorial<0>} wymaga osobnej specjalizacji?
\item Kiedy \texttt{std::decay\_t<T>} i \texttt{std::remove\_cvref\_t<T>}
      dają różne wyniki?
\item Dlaczego \texttt{std::conditional\_t} nie rozwiązuje wszystkich
      problemów warunkowego tworzenia typów?
\item Co dokładnie sprawdza \texttt{void\_t<decltype(++declval<T\&>())>}?
\item Kiedy niepowodzenie podstawienia jest SFINAE, a kiedy zwykłym
      błędem kompilacji?
\item Dlaczego \texttt{if constexpr} nie zastępuje preprocesora?
\item Czym różnią się \texttt{constexpr} i \texttt{consteval}?
\item W jakich sytuacjach CRTP nie zastąpi funkcji wirtualnych?
\item Kiedy lepiej zastosować \texttt{concepts} niż \texttt{enable\_if}?
\end{enumerate}

\paragraph{Dalsza lektura.}
Warto sięgnąć do dokumentacji \emph{cppreference} dotyczącej
\emph{Templates}, \emph{Type traits} i \emph{SFINAE}.
Pomocne są również hasła \texttt{std::void\_t}, \texttt{std::enable\_if},
\texttt{if constexpr}, \texttt{consteval} oraz
\texttt{std::integer\_sequence}. Do diagnozowania złożonych
błędów szablonów przydają się opcje \texttt{-Wall -Wextra -Wpedantic},
\texttt{-ftemplate-backtrace-limit=0} (GCC) oraz ograniczanie zasięgu
instancjowanego kodu do minimalnego przykładu.


\chapter[Metaprogramowanie w C++]{Zaawansowane metaprogramowanie: szablony, cechy typów i obliczenia podczas kompilacji}
\label{chap:meta-cpp}

\section{Czym jest metaprogramowanie?}

\emph{Metaprogramowanie} oznacza pisanie kodu, który na etapie kompilacji
wyznacza typy, własności typów, wartości lub strukturę później
kompilowanego programu. Szablon nie jest gotową funkcją czy klasą:
stanowi przepis, z którego kompilator tworzy konkretne wersje
dla podanych typów. Nazywamy to \emph{instancjowaniem szablonu}.
Na przykład wywołanie
\texttt{twice(21)} w programie poniżej każe kompilatorowi przygotować
wersję funkcji dla \texttt{int}; program podczas działania wywołuje
już tę wersję. W C++20 do dyspozycji mamy kilka narzędzi:
\begin{itemize}
\item \textbf{szablony klas i funkcji} — generowanie kodu dla różnych typów;
\item \textbf{\texttt{type\_traits}} — obliczanie własności i przekształcanie typów;
\item \textbf{specjalizacje i SFINAE} — wybór poprawnych wersji szablonów;
\item \textbf{\texttt{if constexpr}, \texttt{constexpr}, \texttt{consteval}} —
      wykonywanie obliczeń i selekcja kodu w czasie kompilacji;
\item \textbf{\texttt{concepts}} — czytelne formułowanie wymagań wobec typów;
\item \textbf{CRTP} — wykorzystanie dziedziczenia do polimorfizmu statycznego.
\end{itemize}

Metaprogramowanie ma praktyczne cele: eliminowanie powtarzalnego kodu,
uzyskiwanie lepszych błędów kompilacji, wybór wydajnej implementacji oraz
zapewnienie poprawności określonych operacji \emph{przed uruchomieniem} programu.
Nie oznacza automatycznie zwiększenia wydajności — kosztem mogą być dłuższa
kompilacja, powiększenie kodu maszynowego i trudniejsze komunikaty diagnostyczne.

\begin{lstlisting}
#include <type_traits>

template<class T>
T twice(T x) {
    static_assert(std::is_arithmetic_v<T>);
    return x + x;
}

int a = twice(21);       // T = int
// twice("text");       // static_assert fails
\end{lstlisting}

\paragraph{Związki z innymi rozdziałami.}
Wykorzystamy szablony wariadyczne (rozdział~11),
przenoszenie i referencje (rozdział~6) oraz lambdy (rozdział~9).
Koncepcje szerzej poznamy w rozdziale~13; tutaj pokazujemy
przede wszystkim mechanizmy, na których się opierają.
Wszystkie oznaczone jako poprawne przykłady odnoszą się do \textbf{C++20}.

\section{Instancjowanie szablonów, parametry i zależne nazwy}
\label{sec:cpp9-instantiation}

\subsection{Parametry typowe, nietypowe i szablonowe}

Szablon może przyjmować argumenty będące typami, wartościami stałymi
lub innymi szablonami. Parametr wartości (\textit{non-type template parameter},
NTTP) w C++20 może również mieć odpowiedni typ klasowy spełniający
wymagania tzw. typu strukturalnego.

\begin{lstlisting}
#include <array>
#include <cstddef>

template<class T, std::size_t N>
struct Buffer {
    std::array<T, N> data{};
};

Buffer<double, 4> samples; // T=double, N=4
\end{lstlisting}

\texttt{Buffer<double,4>} i \texttt{Buffer<double,8>} są różnymi typami.
Parametr \texttt{N} jest znany podczas kompilacji i może wystąpić tam,
gdzie język wymaga wyrażenia stałego, np. jako rozmiar \texttt{std::array}.

\subsection{Parametr będący obiektem klasy (C++20)}

Od C++20 odpowiednie \emph{typy strukturalne} mogą występować jako
parametry szablonu będące wartościami. Umożliwia to np. identyfikowanie
szablonu przez napis znany podczas kompilacji:

\begin{lstlisting}
#include <cstddef>

template<std::size_t N>
struct fixed_string {
    char data[N]{};
    constexpr fixed_string(const char (&s)[N]) {
        for (std::size_t i = 0; i < N; ++i)
            data[i] = s[i];
    }
};

template<fixed_string Name>
struct NamedCommand {
    static constexpr auto name = Name;
};

static_assert(NamedCommand<"run">::name.data[0] == 'r');
\end{lstlisting}

To przykład zaawansowanej techniki; zwykły obiekt \texttt{std::string}
nie jest w C++20 strukturalnym argumentem szablonu. Rozmiar i
zawartość napisu wpływają tutaj na tożsamość specjalizacji szablonu.

\subsection{Nazwy zależne: \texttt{typename} i \texttt{template}}

W definicji szablonu kompilator nie zawsze wie, czy nazwa zależna od
parametru jest typem, czy wartością. W wielu kontekstach musimy to
jednoznacznie zaznaczyć:

\begin{lstlisting}
#include <vector>

template<class Container>
auto first_copy(const Container& c) {
    typename Container::value_type x = c.front();
    return x;
}

template<class T>
struct Wrapper {
    template<class U>
    static constexpr bool same = false;
};

template<class T>
constexpr bool use_member_template() {
    return Wrapper<T>::template same<int>;
}
\end{lstlisting}

Słowo \texttt{typename} wskazuje, że zależna nazwa członkowska
jest typem. Zapis \texttt{::template same<int>} informuje,
że po operatorze \texttt{::} występuje szablon składowy, a znaki
\texttt{<} i \texttt{>} są delimitatorami argumentów szablonu.
Dokładne reguły zależą od kontekstu; obu słów nie należy dopisywać
mechanicznie wszędzie.

\section{Metafunkcje i \texttt{std::integral\_constant}}
\label{sec:cpp9-integral-constant}

Poniższa silnia obliczana przez szablon pokazuje starszy styl
metaprogramowania. Gdy wynikiem ma być zwykła liczba, funkcja
\texttt{constexpr} jest dziś zazwyczaj prostsza. Szablony klas
pozostają potrzebne zwłaszcza wtedy, gdy wynikiem obliczenia ma być
\emph{typ}.

Klasyczne metaprogramowanie szablonowe opierało się na klasach, które
\emph{przenosiły wyniki obliczenia w swoich składowych statycznych}.
Dla przykładu silnię można wyznaczyć przez rekurencyjne instancjowanie
szablonów:

\begin{lstlisting}
template<unsigned N>
struct Factorial {
    static constexpr unsigned value =
        N * Factorial<N - 1>::value;
};

template<>
struct Factorial<0> {
    static constexpr unsigned value = 1;
};

static_assert(Factorial<5>::value == 120);
\end{lstlisting}

Specjalizacja dla zera pełni rolę warunku stopu. Bez niej kompilator
próbowałby tworzyć kolejne instancje, aż osiągnąłby limit głębokości.
Obliczenie ma charakter rekursji \emph{w strukturze typów}, a nie
rekursji wywołań funkcji wykonywanych w czasie działania programu.

Standardowym fundamentem takich metafunkcji jest
\texttt{std::integral\_constant<T,v>}:

\begin{lstlisting}
#include <type_traits>

using TrueType = std::integral_constant<bool, true>;
static_assert(TrueType::value);
static_assert(std::bool_constant<true>::value);
static_assert(std::true_type::value);
static_assert(!std::false_type::value);

template<unsigned N>
struct PowerOfTwo :
    std::integral_constant<unsigned, 2 * PowerOfTwo<N-1>::value> {};

template<>
struct PowerOfTwo<0> : std::integral_constant<unsigned, 1> {};

static_assert(PowerOfTwo<5>::value == 32);
\end{lstlisting}

Predykaty takie jak \texttt{std::is\_pointer} udostępniają wartość
\texttt{::value}. Od C++17 można używać krótszej wersji
z przyrostkiem \texttt{\_v}, na przykład:
\begin{center}
\texttt{std::is\_pointer\_v<T>}.
\end{center}

\subsection{Metafunkcja szablonowa czy funkcja \texttt{constexpr}?}

\begin{lstlisting}
constexpr unsigned factorial(unsigned n) {
    unsigned result = 1;
    for (unsigned i = 2; i <= n; ++i)
        result *= i;
    return result;
}

static_assert(factorial(5) == 120);
unsigned x = 6;
unsigned y = factorial(x); // allowed at run time
\end{lstlisting}

Jeśli wynik jest zwykłą wartością, funkcja \texttt{constexpr} jest
najczęściej prostsza niż rekurencyjna metafunkcja. Klasyczne metafunkcje
pozostają użyteczne, gdy \emph{wynikiem jest typ}, gdy korzystamy ze
specjalizacji lub interoperujemy ze starszym API.

\section{Specjalizacje: pełne i częściowe}
\label{sec:cpp9-specialization}

Szablon podstawowy definiuje przypadek ogólny. \emph{Specjalizacja pełna}
dotyczy konkretnego zestawu argumentów, a \emph{specjalizacja częściowa}
obsługuje pewną ich rodzinę:

\begin{lstlisting}
#include <type_traits>

template<class T>
struct is_raw_pointer : std::false_type {};

template<class T>
struct is_raw_pointer<T*> : std::true_type {};

static_assert(!is_raw_pointer<int>::value);
static_assert(is_raw_pointer<double*>::value);
static_assert(!is_raw_pointer<int* const>::value);
\end{lstlisting}

Ostatni wynik nie jest pomyłką: \texttt{int* const} jest
\emph{stałym wskaźnikiem}, nie typem postaci \texttt{T*}.
Jeśli chcemy ignorować kwalifikatory zewnętrzne, możemy najpierw
użyć \texttt{std::remove\_cv\_t<T>}.

\begin{lstlisting}
template<class T>
inline constexpr bool is_any_raw_pointer_v =
    is_raw_pointer<std::remove_cv_t<T>>::value;

static_assert(is_any_raw_pointer_v<int* const>);
\end{lstlisting}

\paragraph{Ważne ograniczenie.}
Dla \textbf{szablonów funkcji nie ma specjalizacji częściowej}.
W takich przypadkach wykorzystujemy przeciążanie, koncepcje,
\texttt{if constexpr} lub specjalizację pomocniczego szablonu klasy.
Pełna specjalizacja funkcji istnieje, ale jej interakcja z przeciążeniami
bywa zaskakująca; w nowym kodzie zwykle preferujemy przeciążenia.

\section{\texttt{type\_traits}: pytania o typy}
\label{sec:cpp9-traits}

Biblioteka \texttt{<type\_traits>} jest zbiorem standardowych metafunkcji.
Dzielimy je na \emph{predykaty} (wynik logiczny) i \emph{transformacje}
(wynik będący typem). Ich składnia opiera się na ustalonych konwencjach:
\begin{itemize}
\item \texttt{Trait<T>::value} lub \texttt{Trait\_v<T>} dla predykatów;
\item \texttt{typename Trait<T>::type} lub \texttt{Trait\_t<T>}
      dla transformacji.
\end{itemize}

\begin{lstlisting}
#include <type_traits>
#include <string>

static_assert(std::is_integral_v<int>);
static_assert(std::is_floating_point_v<double>);
static_assert(std::is_same_v<int, int>);
static_assert(std::is_pointer_v<const char*>);
static_assert(std::is_constructible_v<std::string, const char*>);
static_assert(std::is_move_constructible_v<std::string>);

using A = std::remove_const_t<const int>; // int
using B = std::remove_reference_t<int&>; // int
using C = std::remove_cvref_t<const int&>; // int (C++20)
using D = std::add_pointer_t<int>;        // int*
static_assert(std::is_same_v<A, B>);
static_assert(std::is_same_v<C, int>);
static_assert(std::is_same_v<D, int*>);
\end{lstlisting}

\begin{center}
\small
\begin{tabular}{|>{\raggedright\arraybackslash}p{5.5cm}|>{\raggedright\arraybackslash}p{7.0cm}|}
\hline
\textbf{Cecha} & \textbf{Co sprawdza lub wyznacza?} \\
\hline
\texttt{is\_same\_v<A,B>} & Identyczność typów (wraz z istotnymi kwalifikatorami). \\
\hline
\texttt{is\_convertible\_v<A,B>} & Dostępność konwersji niejawnej \texttt{A} do \texttt{B}. \\
\hline
\texttt{is\_constructible\_v<T,Args...>} & Możliwość skonstruowania \texttt{T} z argumentów zadanych typów. \\
\hline
\texttt{is\_nothrow\_move\_}\newline\texttt{constructible\_v<T>} & Deklarowana bezwyjątkowość konstrukcji przenoszącej. \\
\hline
\texttt{is\_trivially\_copyable\_v<T>} & Własność pozwalająca na określone operacje z reprezentacją bajtową. \\
\hline
\texttt{remove\_cvref\_t<T>} & Usunięcie referencji i kwalifikatorów \texttt{const}/\texttt{volatile} z poziomu wierzchniego. \\
\hline
\texttt{common\_type\_t<A,B>} & Wspólny typ wynikowy dla danych typów (o ile istnieje). \\
\hline
\texttt{invoke\_result\_t<F,Args...>} & Typ wyniku wywołania poprzez \texttt{std::invoke}. \\
\hline
\end{tabular}
\end{center}

\paragraph{\texttt{remove\_cvref} a \texttt{decay}.}
\texttt{std::decay\_t<T>} oprócz usunięcia referencji i zewnętrznych
kwalifikatorów może przekształcić tablicę w wskaźnik oraz typ funkcji
w wskaźnik do funkcji. Nie są to zatem synonimy.

\begin{lstlisting}
using X = std::remove_cvref_t<int(&)[3]>; // int[3]
using Y = std::decay_t<int(&)[3]>;        // int*
static_assert(std::is_same_v<Y, int*>);
\end{lstlisting}

\paragraph{Czego predykat nie gwarantuje?}
\texttt{std::is\_move\_constructible\_v<T>} może być prawdziwe również
wtedy, gdy wywołanie z argumentem rvalue używa konstruktora kopiującego.
Podobnie \texttt{is\_trivially\_copyable} nie oznacza, że każdy odczyt
bajtów, rzutowanie lub serializacja do pliku będzie poprawna bądź
przenośna. Cechy typów opisują wybrane własności języka, nie
pełny kontrakt semantyczny.

\subsection{Składanie predykatów z zachowaniem krótkiego obliczania}

\texttt{std::conjunction}, \texttt{std::disjunction} i
\texttt{std::negation} łączą predykaty, które mają składową
\texttt{::value}. Dwa pierwsze narzędzia zatrzymują odczytywanie \texttt{::value}
z kolejnych predykatów, kiedy wynik jest już znany. Trzeba jednak
poprawnie sformować same argumenty typowe tych metafunkcji.

\begin{lstlisting}
#include <type_traits>

template<class T>
using SmallSignedInteger = std::conjunction<
    std::is_integral<T>,
    std::is_signed<T>,
    std::bool_constant<(sizeof(T) <= sizeof(int))>>;

static_assert(SmallSignedInteger<short>::value);
static_assert(!SmallSignedInteger<unsigned>::value);
\end{lstlisting}

Dla prostych stałych logicznych zapis
\texttt{(cond1 \&\& cond2)} bywa czytelniejszy.
\texttt{conjunction} jest użyteczne przede wszystkim wtedy, gdy
poszczególne warunki są odrębnymi metafunkcjami, a ich instancjowanie
może być kosztowne lub niepoprawne bez wcześniejszego warunku.

\section{Metafunkcje zwracające typy i leniwy wybór}
\label{sec:cpp9-transforms}

Tworząc własną transformację typu, zwyczajowo zapisujemy wynik jako
\texttt{using type = ...} i udostępniamy wygodny alias:

\begin{lstlisting}
template<class T>
struct storage_type {
    using type = std::remove_cvref_t<T>;
};

template<class T>
using storage_type_t = typename storage_type<T>::type;

static_assert(std::is_same_v<storage_type_t<const long&>, long>);
\end{lstlisting}

Ważnym narzędziem jest \texttt{std::conditional\_t<B,T,F>}, wybierające
jeden z dwóch typów w zależności od znanej w kompilacji wartości logicznej:

\begin{lstlisting}
template<class T>
using accumulator_t = std::conditional_t<
    std::is_integral_v<T>, long long, long double>;

static_assert(std::is_same_v<accumulator_t<int>, long long>);
static_assert(std::is_same_v<accumulator_t<float>, long double>);
\end{lstlisting}

\paragraph{Pułapka eager instantiation.}
\texttt{conditional\_t} wybiera typ spośród \emph{już utworzonych
argumentów typowych}. Nie chroni przed błędami powstającymi podczas
formowania któregoś z nich:

\begin{lstlisting}
// Incorrect for int: typename T::value_type is invalid.
// template<class T>
// using bad = std::conditional_t<
//     std::is_integral_v<T>, T, typename T::value_type>;
\end{lstlisting}

Aby uzyskać prawdziwie warunkowe tworzenie typu, można zastosować
specjalizację pomocniczej klasy:

\begin{lstlisting}
template<class T, bool = std::is_integral_v<T>>
struct value_or_self {
    using type = typename T::value_type;
};

template<class T>
struct value_or_self<T, true> {
    using type = T;
};

template<class T>
using value_or_self_t = typename value_or_self<T>::type;

static_assert(std::is_same_v<value_or_self_t<int>, int>);
\end{lstlisting}

Podstawowa definicja zawiera nazwę zależną, która jest sprawdzana przy
instancjowaniu odpowiedniego przypadku; nie jest potrzebna dla \texttt{int}.

\section{\texttt{decltype}, \texttt{declval} i konteksty bez obliczania}

\begin{center}\small
\begin{tabularx}{\textwidth}{@{}p{0.32\textwidth}X@{}}
\toprule
Narzędzie & Za co odpowiada \\
\midrule
\code{decltype(expr)} & Ustala typ wyrażenia bez wykonywania go; szczególna reguła dotyczy nieujętej w nawiasy nazwy. \\
\code{decltype((x))} & Dla lvalue \code{x} zwraca typ referencyjny, np. \texttt{T\&}. \\
\code{std::declval<T>()} & Udostępnia fikcyjne wyrażenie typu \code{T} wyłącznie w kontekście nieobliczanym. \\
\code{std::void_t<...>} & Zamienia poprawne podstawienia typów na \code{void}; przydatne w starszych technikach wykrywania. \\
\code{if constexpr} & W szablonie wybiera gałąź podczas kompilacji i odrzuca niepotrzebną gałąź po instancjacji. \\
\bottomrule
\end{tabularx}
\end{center}
\label{sec:cpp9-decltype}

\texttt{decltype(expr)} pozwala odczytać typ wyrażenia bez wykonywania
samego wyrażenia. Subtelność stanowi różnica między nieujętym w nawiasy
identyfikatorem a ogólnym wyrażeniem:

\begin{lstlisting}
int x = 5;
static_assert(std::is_same_v<decltype(x), int>);
static_assert(std::is_same_v<decltype((x)), int&>);
\end{lstlisting}

Wyrażenie \texttt{(x)} jest \textit{lvalue}, więc \texttt{decltype((x))}
zwraca \texttt{int\&}. Reguła ta ma bezpośredni związek z kategoriami
wartości (rozdział~6).

\texttt{std::declval<T>()} pozwala użyć argumentu typu \texttt{T}
w \emph{kontekście niewartościowanym}, nawet jeśli nie możemy utworzyć
jego obiektu. Nie wolno rzeczywiście wykonywać wywołania
\texttt{declval}.

\begin{lstlisting}
#include <utility>
#include <type_traits>

template<class T>
using dereference_t = decltype(*std::declval<T&>());

static_assert(std::is_same_v<dereference_t<int*>, int&>);

struct NonDefault {
    NonDefault() = delete;
    int method() const;
};

using Result = decltype(std::declval<const NonDefault&>().method());
static_assert(std::is_same_v<Result, int>);
\end{lstlisting}

Typ wyniku wywołania można również wyznaczyć za pomocą gotowej metafunkcji:
\begin{center}
\texttt{std::invoke\_result\_t<F, Args...>}.
\end{center}

\section{SFINAE: błędne podstawienie to nie zawsze błąd kompilacji}
\label{sec:cpp9-sfinae}

Ta część pomaga czytać biblioteki napisane przed C++20.
W nowym kodzie wymagania wobec typów zwykle prościej wyrazić
przez \texttt{requires}, opisane w rozdziale~13.

\textbf{SFINAE} (\textit{Substitution Failure Is Not An Error}) mówi,
że niektóre błędy powstałe podczas podstawiania argumentów do szablonu
\emph{nie kończą kompilacji}: dana specjalizacja zostaje wyłączona
z zestawu kandydatów lub odpowiednia specjalizacja częściowa nie pasuje.
Dotyczy to określonych kontekstów bezpośredniego podstawiania;
\textbf{nie} jest mechanizmem przechwytującym dowolne błędy w ciele funkcji
ani skutki uboczne instancjowania innych typów.

\subsection{\texttt{std::enable\_if\_t}}

\begin{lstlisting}
#include <type_traits>

template<class T>
constexpr std::enable_if_t<std::is_integral_v<T>, T>
doubled(T x) {
    return x + x;
}

static_assert(doubled(3) == 6);
// doubled(2.5); // no matching function
\end{lstlisting}

Dla typu \texttt{double} nie istnieje \texttt{type} w
\texttt{std::enable\_if<false,double>}; funkcja zostaje wyłączona
z odpowiedniego zbioru kandydatów. Do konstruktorów, operatorów i
różnych wariantów przeciążeń istnieją odmienne miejsca umieszczania
\texttt{enable\_if}; łatwo przy tym uzyskać przypadkowe powtórne
deklaracje lub nieczytelne błędy.

\subsection{\texttt{void\_t} — detection idiom}

Od C++17 wygodnie wykrywać istnienie nazw i poprawność wyrażeń przy
użyciu \texttt{std::void\_t}. Jeżeli wszystkie podane argumenty
są prawidłowymi typami, rezultatem jest po prostu \texttt{void}:

\begin{lstlisting}
#include <type_traits>
#include <vector>

template<class, class = void>
struct has_value_type : std::false_type {};

template<class T>
struct has_value_type<T, std::void_t<typename T::value_type>>
    : std::true_type {};

static_assert(has_value_type<std::vector<int>>::value);
static_assert(!has_value_type<int>::value);
\end{lstlisting}

W podstawieniu \texttt{T=int} wyrażenie \texttt{int::value\_type}
nie ma sensu; specjalizacja częściowa nie jest więc wybierana.
Analogicznie można badać poprawność wyrażenia z \texttt{decltype}:

\begin{lstlisting}
#include <type_traits>
#include <utility>

template<class, class = void>
struct has_preincrement : std::false_type {};

template<class T>
struct has_preincrement<T,
    std::void_t<decltype(++std::declval<T&>())>>
    : std::true_type {};

static_assert(has_preincrement<int>::value);
\end{lstlisting}

\subsection{SFINAE kontra \texttt{requires} (C++20)}

Konstrukcję wykrywającą operację możemy wyrazić bezpośrednio za pomocą
koncepcji z rozdziału~13:

\begin{lstlisting}
#include <concepts>

template<class T>
concept Preincrementable = requires(T& x) {
    { ++x } -> std::same_as<T&>;
};

template<Preincrementable T>
void advance_one(T& x) { ++x; }
\end{lstlisting}

Metafunkcja \texttt{has\_preincrement<T>} sprawdza jedynie poprawność
operacji \texttt{++x}. Z kolei koncepcja \texttt{Preincrementable}
dodatkowo wymaga, aby jej \emph{typ wyniku} odpowiadał dokładnie \texttt{T\&}.
Różnica ta jest zamierzona: ograniczenia należy definiować zgodnie z
kontraktem, a nie tylko według dostępności pewnego wyrażenia.

\section{\texttt{if constexpr}: wybór implementacji na podstawie typu}
\label{sec:cpp9-ifconstexpr}

Od C++17 często zamiast rekurencyjnych specjalizacji wybieramy gałąź
wewnątrz pojedynczego szablonu. Niebrana gałąź zależna od argumentu
szablonu nie jest instancjowana dla danego zestawu argumentów:

\begin{lstlisting}
#include <type_traits>

template<class T>
constexpr auto number_or_value(const T& x) {
    if constexpr (std::is_arithmetic_v<T>)
        return x;
    else
        return x.value();
}

struct Box {
    constexpr int value() const { return 7; }
};
static_assert(number_or_value(3) == 3);
static_assert(number_or_value(Box{}) == 7);
\end{lstlisting}

\paragraph{Dlaczego nie zwykłe \texttt{if}?}
W zwykłym \texttt{if} kompilator musi przeanalizować obie gałęzie
pod kątem poprawności dla instancjowanego kodu. Dla \texttt{int}
wyrażenie \texttt{x.value()} nie jest poprawne, a sam warunek
\texttt{std::is\_arithmetic\_v<int>} tego nie zmieni.

\paragraph{Ważna subtelność.}
\texttt{if constexpr(false)} nie jest zamiennikiem
\texttt{\#if 0}: poza kontekstem szablonowym odrzucona gałąź nadal
podlega standardowym kontrolom poprawności. Ponadto mechanizm nie
usprawiedliwia konstrukcji zawsze niepoprawnych niezależnie od
argumentów szablonu.

\section{Pakiety typów, sekwencje indeksów i fold expressions}
\label{sec:cpp9-type-lists}

Metaprogramowanie dobrze współpracuje z pakietami parametrów
(rozdział~11). Zamiast tworzyć klasyczne rekurencyjne listy typów,
często możemy używać rozwinięcia pakietu i wyrażenia zwijającego.

\begin{lstlisting}
#include <type_traits>

template<class... Ts>
struct type_list {};

template<class T, class... Ts>
inline constexpr bool contains_v =
    (std::is_same_v<T, Ts> || ...); // empty pack -> false

static_assert(contains_v<int, char, int, double>);
static_assert(!contains_v<long, char, int, double>);
static_assert(!contains_v<int>);
\end{lstlisting}

Zamiast tworzyć specjalizacje dla każdego elementu pakietu, korzystamy
z jednego \texttt{fold expression}. Gdy potrzebna jest reprezentacja
indeksów \texttt{0,1,...,N-1}, biblioteka udostępnia
\texttt{std::make\_index\_sequence<N>}:

\begin{lstlisting}
#include <tuple>
#include <utility>
#include <cstddef>

template<class Tuple, std::size_t... I>
constexpr auto reversed_impl(const Tuple& t,
                             std::index_sequence<I...>) {
    return std::make_tuple(std::get<sizeof...(I)-1-I>(t)...);
}

template<class... Ts>
constexpr auto reversed(const std::tuple<Ts...>& t) {
    return reversed_impl(t, std::index_sequence_for<Ts...>{});
}

static_assert(reversed(std::tuple{1, 2, 3}) ==
              std::tuple{3, 2, 1});
\end{lstlisting}

\paragraph{Czy zawsze potrzebujemy indeksów?}
Nie. W przypadku prostego wykonania funkcji na elementach krotki
w C++17 często wystarczy \texttt{std::apply}; sekwencja indeksów
jest potrzebna zwłaszcza wtedy, gdy chcemy \emph{przekształcić pozycje}
lub mieć indywidualny dostęp do każdego indeksu na etapie kompilacji.

\section{\texttt{constexpr}, \texttt{consteval} i obliczenia na danych}
\label{sec:cpp9-constant-eval}

Szablony nie są jedynym sposobem wykonania pracy w trakcie kompilacji.
Od C++11 rozwijano możliwości \texttt{constexpr}, a w C++20 mamy
również \texttt{consteval}. Różnica jest istotna:
\begin{itemize}
\item funkcja \texttt{constexpr} \emph{może} uczestniczyć w obliczeniach
stałych, a w dozwolonych wywołaniach działa też podczas wykonania;
\item wywołanie funkcji \texttt{consteval} musi spełniać reguły
właściwe obliczaniu w czasie kompilacji (\textit{immediate function}).
\end{itemize}

\begin{lstlisting}
constexpr int square(int x) { return x * x; }
consteval int cube_ct(int x) { return x * x * x; }

static_assert(square(4) == 16);
static_assert(cube_ct(3) == 27);

int n = 4;
int a = square(n); // runtime evaluation possible
// int b = cube_ct(n); // error: n is not constant
\end{lstlisting}

Przykładem praktycznym jest haszowanie \texttt{std::string\_view}
w stałym wyrażeniu. Poniższy wariant FNV-1a nadaje się do zbudowania
stałej identyfikującej np. nazwę polecenia. Nie jest to hash
kryptograficzny i nie wyklucza kolizji.

\begin{lstlisting}
#include <cstdint>
#include <string_view>

consteval std::uint64_t fnv1a(std::string_view text) {
    std::uint64_t h = 14695981039346656037ull;
    for (unsigned char c : text) {
        h ^= c;
        h *= 1099511628211ull;
    }
    return h;
}

static_assert(fnv1a("abc") == 0xe71fa2190541574bull);
\end{lstlisting}

\paragraph{\texttt{constinit} i \texttt{std::is\_constant\_evaluated}.}
\texttt{constinit} zapewnia statyczną inicjalizację odpowiedniej
zmiennej, lecz sam nie czyni jej stałą. Funkcja
\texttt{std::is\_constant\_evaluated()} (C++20) pozwala w
pewnych przypadkach rozpoznać, czy bieżąca ewaluacja ma charakter
stały; nadużywanie jej może jednak zaciemniać semantykę funkcji.

\section{CRTP: polimorfizm statyczny}
\label{sec:cpp9-crtp}

\textbf{CRTP} (\textit{Curiously Recurring Template Pattern}) polega na
tym, że klasa dziedziczy po specjalizacji szablonu, której argumentem
jest ona sama:

\begin{lstlisting}
template<class Derived>
struct Printable {
    void print() const {
        static_cast<const Derived&>(*this).print_impl();
    }
};

struct Number : Printable<Number> {
    int value = 0;
    void print_impl() const;
};
\end{lstlisting}

Mechanizm \texttt{static\_cast} jest tutaj poprawny pod założeniem,
że obiekt bazowy faktycznie stanowi podobiekt \texttt{Derived}.
CRTP nie zapewnia dynamicznego sprawdzania typu ani wirtualnego
wywołania — poprawność wynika z konstrukcji hierarchii. Wywołanie
\texttt{print\_impl()} może być statycznie związane i często
łatwiej je wbudować w kod wynikowy (\textit{inline}).

Bardziej praktycznym przykładem jest domieszka (\textit{mixin})
dostarczająca wspólną operację, lecz działająca na danych
konkretnej klasy pochodnej:

\begin{lstlisting}
template<class Derived>
struct TranslationMixin {
    constexpr void translate(int dx, int dy) {
        auto& self = static_cast<Derived&>(*this);
        self.x += dx;
        self.y += dy;
    }
};

struct Point : TranslationMixin<Point> {
    int x{}, y{};
    constexpr Point(int a, int b) : x(a), y(b) {}
};

constexpr bool test_translation() {
    Point p{1, 2};
    p.translate(3, 4);
    return p.x == 4 && p.y == 6;
}
static_assert(test_translation());
\end{lstlisting}

Baza \texttt{TranslationMixin<Point>} jest instancjowana z typem
\texttt{Point}, chociaż sam typ \texttt{Point} nie jest jeszcze
kompletny w momencie deklarowania dziedziczenia. Odwołania do składowych
\texttt{x} i \texttt{y} zależą jednak od \texttt{Derived} i są
sprawdzane dopiero podczas instancjowania odpowiedniej funkcji.
W rzeczywistym API warto opisać wymagane składowe koncepcją lub
przenieść operację do jawnie wskazanego interfejsu.

\paragraph{CRTP a wirtualne funkcje.}
Klasyczny polimorfizm dynamiczny pozwala wybierać implementację
na podstawie rzeczywistego typu obiektu przez interfejs klasy bazowej,
także gdy typ jest nieznany podczas kompilacji kodu wywołującego.
CRTP wymaga znajomości konkretnego typu podczas instancjowania.
Nie jest ogólnym zamiennikiem klas abstrakcyjnych. W C++20 prostsze
zastosowania domieszek i interfejsów statycznych można często oprzeć
na \texttt{concepts} i funkcjach wolnostojących.

\section{Policy-based design: wybór zachowania przez typ}
\label{sec:cpp9-policy}

Szablon może przyjąć \emph{politykę}, której metody definiują wariant
algorytmu. Umożliwia to wybór zachowania bez wirtualnych wywołań,
a kompilator może całkowicie usunąć niepotrzebne operacje:

\begin{lstlisting}
#include <limits>
#include <stdexcept>

struct Checked {
    static int divide(int a, int b) {
        if (b == 0) throw std::invalid_argument("zero");
        if (a == std::numeric_limits<int>::min() && b == -1)
            throw std::overflow_error("result outside int");
        return a / b;
    }
};

struct Fast {
    static int divide(int a, int b) { return a / b; }
};

template<class Policy>
struct Calculator {
    int divide(int a, int b) const {
        return Policy::divide(a, b);
    }
};

// Calculator<Checked> checks both invalid integer cases;
// Calculator<Fast> requires a nonzero divisor and no overflow.
\end{lstlisting}

Polityka \texttt{Fast} nie zabezpiecza dzielenia przez zero ani
przepełnienia \texttt{INT\_MIN / -1}; oba przypadki mają
niezdefiniowane zachowanie dla \texttt{int}. \texttt{Checked}
odrzuca je przed dzieleniem. Dobór polityki jest częścią kontraktu API.

W realnym kodzie szablon polityki można ograniczyć koncepcją, a jeśli
liczba wariantów jest niewielka, zwykłe \texttt{if constexpr} może być
łatwiejsze w utrzymaniu niż rozbudowana hierarchia polityk.

\section{Kompletny przykład: dobór typu akumulatora}
\label{sec:cpp9-worked}

Załóżmy, że chcemy obliczać sumę liczb, lecz wybierać typ akumulatora
w zależności od argumentu: dla typów całkowitych \texttt{long long},
a dla zmiennoprzecinkowych \texttt{long double}. Rozwiązanie łączy
\texttt{concepts}, \texttt{type\_traits} i \texttt{conditional\_t}:

\begin{lstlisting}
#include <concepts>
#include <type_traits>
#include <vector>

template<class T>
concept Numeric = std::integral<T> || std::floating_point<T>;

template<Numeric T>
using sum_type_t = std::conditional_t<
    std::integral<T>, long long, long double>;

template<Numeric T>
sum_type_t<T> accumulate_values(const std::vector<T>& xs) {
    sum_type_t<T> result{};
    for (const T& x : xs) {
        result += static_cast<sum_type_t<T>>(x);
    }
    return result;
}

static_assert(std::is_same_v<sum_type_t<int>, long long>);
static_assert(std::is_same_v<sum_type_t<double>, long double>);
\end{lstlisting}

\paragraph{Ograniczenia.}
Przykład wyjaśnia selekcję typu, lecz \emph{nie daje gwarancji
braku przepełnienia}. Nawet \texttt{long long} ma ograniczony zakres,
a suma zmiennoprzecinkowa może akumulować błędy zaokrągleń.
Najpierw należy ustalić wymagania numeryczne, a dopiero potem
projektować odpowiednią metafunkcję.

\section{Czas kompilacji, diagnostyka i praktyczne ograniczenia}

Metaprogramowanie przenosi koszt części pracy z czasu wykonania na
czas kompilacji, ale samo w sobie nie oznacza wydajniejszego programu.
Każda odmienna specjalizacja szablonu może wymagać osobnego
instancjowania i wygenerowania kodu. Powstają wtedy m.in.:
\begin{itemize}
\item dłuższe kompilacje i większe zużycie pamięci kompilatora;
\item \emph{code bloat}, czyli powiększenie plików wynikowych przez wiele
podobnych instancji;
\item skomplikowane komunikaty błędów przy zagnieżdżonych metafunkcjach;
\item problemy z zależnościami między plikami nagłówkowymi;
\item ograniczenia głębokości instancjowania i kosztu ewaluacji stałej.
\end{itemize}

\paragraph{Pliki nagłówkowe i jawne instancjowanie.}
Definicje szablonów zazwyczaj umieszczamy w nagłówkach, ponieważ
kompilator musi je znać przy instancjowaniu z konkretnymi argumentami.
Możliwe jest również \emph{jawne instancjowanie} dla ustalonych typów,
co pozwala kontrolować zestaw generowanych specjalizacji:

\begin{lstlisting}
// In a header:
template<class T> T triple(T x);
extern template int triple<int>(int);

// In one .cpp file (definition must be visible here):
template<class T> T triple(T x) { return x + x + x; }
template int triple<int>(int);
\end{lstlisting}

Ten przykład przedstawia ideę rozdzielenia deklaracji i jawnego
instancjowania; obie oznaczone części należy umieścić w odpowiednich
plikach, a dla wywołań innych typów zapewnić dostęp do definicji.
Warto pamiętać o regule ODR z rozdziału~4.

\section{Jaką technikę wybrać?}

\begin{center}
\small
\begin{tabular}{|p{5.7cm}|p{6.8cm}|}
\hline
\textbf{Potrzeba} & \textbf{Najczęściej odpowiedni mechanizm} \\
\hline
Sprawdzenie wymaganych operacji w nowym API & \texttt{concepts}, \texttt{requires} (C++20) \\
\hline
Obliczenie wartości dla stałych danych & Funkcja \texttt{constexpr} \\
\hline
Wymuszenie ewaluacji w czasie kompilacji & \texttt{consteval} (C++20) \\
\hline
Warunkowy kod zależny od typu & \texttt{if constexpr} \\
\hline
Ustalenie własności typu & \texttt{std::is\_...\_v}, inne \texttt{type\_traits} \\
\hline
Zmiana lub wybór typu wynikowego & \texttt{...\_t}, specjalizacja pomocniczej klasy \\
\hline
Wsparcie biblioteki sprzed C++20 & SFINAE, \texttt{enable\_if}, \texttt{void\_t} \\
\hline
Wieloargumentowe cechy typów & Pakiety typów, \texttt{fold expressions} \\
\hline
Dostęp do pozycji elementów krotki & \texttt{index\_sequence}, \texttt{std::apply} \\
\hline
Wielokrotne użycie zachowania klasowego & CRTP / mixin, rozważ też koncepcje \\
\hline
Wybór wariantu strategii w typie & Szablon polityki (policy-based design) \\
\hline
\end{tabular}
\end{center}

\section{Typowe błędy i zasady projektowe}

\begin{enumerate}
\item \textbf{Mylenie typu z wartością.} \texttt{typename T::value\_type}
      to typ, a \texttt{Trait<T>::value} to wartość; nie są wymienne.
\item \textbf{Brak warunku stopu.} Rekurencyjny szablon klasy może
      prowadzić do nieograniczonego instancjowania i błędów głębokości.
\item \textbf{Założenie, że SFINAE przechwytuje wszystkie błędy.}
      Działa tylko w określonych kontekstach podstawiania.
\item \textbf{Nadużywanie \texttt{conditional\_t}.} Oba argumenty
      typowe trzeba poprawnie sformować, zanim zostanie dokonany wybór.
\item \textbf{Próba wywołania \texttt{declval}.} Można go używać wyłącznie
      w odpowiednich kontekstach niewartościowanych.
\item \textbf{Mylenie \texttt{decltype(x)} z \texttt{decltype((x))}.}
      Dodatkowe nawiasy mogą zmienić wynik na typ referencyjny.
\item \textbf{Mylenie \texttt{constexpr} z obowiązkiem obliczenia
      podczas kompilacji.} Nie każde wywołanie funkcji \texttt{constexpr}
      jest ewaluowane w czasie kompilacji.
\item \textbf{Przenoszenie semantyki do traits.} Typ spełniający
      syntaktyczne wymaganie nie musi spełniać własności biznesowej,
      np. braku utraty precyzji czy bezpiecznej serializacji.
\item \textbf{Niepotrzebny CRTP.} Gdy wystarczy wolna funkcja
      z koncepcją, hierarchia klas często niczego nie upraszcza.
\item \textbf{Przedwczesne metaprogramowanie.} Najprostsze czytelne
      rozwiązanie jest zazwyczaj lepsze od dodatkowej warstwy abstrakcji.
\end{enumerate}

\paragraph{Pytania kontrolne.}
\begin{enumerate}
\item Czym różni się instancjowanie szablonu od wykonania funkcji?
\item Dlaczego \texttt{Factorial<0>} wymaga osobnej specjalizacji?
\item Kiedy \texttt{std::decay\_t<T>} i \texttt{std::remove\_cvref\_t<T>}
      dają różne wyniki?
\item Dlaczego \texttt{std::conditional\_t} nie rozwiązuje wszystkich
      problemów warunkowego tworzenia typów?
\item Co dokładnie sprawdza \texttt{void\_t<decltype(++declval<T\&>())>}?
\item Kiedy niepowodzenie podstawienia jest SFINAE, a kiedy zwykłym
      błędem kompilacji?
\item Dlaczego \texttt{if constexpr} nie zastępuje preprocesora?
\item Czym różnią się \texttt{constexpr} i \texttt{consteval}?
\item W jakich sytuacjach CRTP nie zastąpi funkcji wirtualnych?
\item Kiedy lepiej zastosować \texttt{concepts} niż \texttt{enable\_if}?
\end{enumerate}

\paragraph{Dalsza lektura.}
Warto sięgnąć do dokumentacji \emph{cppreference} dotyczącej
\emph{Templates}, \emph{Type traits} i \emph{SFINAE}.
Pomocne są również hasła \texttt{std::void\_t}, \texttt{std::enable\_if},
\texttt{if constexpr}, \texttt{consteval} oraz
\texttt{std::integer\_sequence}. Do diagnozowania złożonych
błędów szablonów przydają się opcje \texttt{-Wall -Wextra -Wpedantic},
\texttt{-ftemplate-backtrace-limit=0} (GCC) oraz ograniczanie zasięgu
instancjowanego kodu do minimalnego przykładu.
